\subsection{局部环上的 Macaulay 模}

命题 3.--- 设 A 为 Noether 局部环，M 为非零的有限生成 A-模，d 为 M 的维数。下列条件等价：

\begin{enumerate}
\def\labelenumi{(\roman{enumi})}
\item
  A-模 M 是 Macaulay 的
\item
  有 \(\operatorname{prof}(\mathbf{M}) = d\)；
\item
  有 \(\mathrm{Ext}_{\Lambda}^{i}(\kappa_{\mathrm{A}},\mathrm{M}) = 0\)（对每个整数 \(i < d\)）；
\item
  有 \(\mathrm{Ext}_{\mathrm{A}}^{i}(\mathrm{N},\mathrm{M}) = 0\)（对每个有限长度的 A-模 N 及每个整数 \(i < d\)）；
\item
  有 \(\operatorname{Ext}_A^i (\mathrm{N},\mathrm{M}) = 0\)（对每个有限生成 A-模 N 及每个整数 \(i < d - \dim (\mathrm{M}\otimes_{\mathrm{A}}\mathrm{N})\)）；
\item
  存在一个由 \(\mathfrak{m}_{\mathrm{A}}\) 中元素组成的长度为 \(d\) 的 M-正则序列。
\end{enumerate}

条件 (i) 等价于等式 prof(M) = d，亦即等价于条件 (ii)，或者再等价于不等式 prof(M) ≥ d，也就是等价于 (iii) 与 (vi) (§ 1, 第 4 节, 定理 2)。蕴含关系 (v) ⇒ (iv) 与 (iv) ⇒ (iii) 是显然的。

最后，设 M 是 Macaulay 的，并设 N 为一个有限生成 A-模。令 \(F = \text{Supp}(N)\)；由 II, § 4, 第 4 节, 命题 18，有 \(\text{Supp}(M) \cap F = \text{Supp}(M \otimes_{A} N)\) ，从而

\[
\operatorname{prof} _ {F} (M) = \operatorname{codim} (\operatorname{Supp} (M) \cap F, \operatorname{Supp} (M)) = \dim M - \dim (M \otimes_ {A} N)
\]

(第 1 节, 命题 1 的推论及第 2 节, 命题 2 的推论)。蕴含关系 (i)⇒(v) 于是由 § 1, 第 5 节, 注 1 得出。

在本节其余部分，我们称 Noether 局部环 A 上的有限生成模 M 是纯的，如果对 M 的每个伴随素理想 p，都有 \(\dim(\mathrm{A}/\mathfrak{p}) = \dim(\mathrm{M})\)。这也意味着 M 没有嵌入伴随素理想，并且 M 的支撑集的诸不可约分支都具有相同的维数。Noether 局部环上的每个 Macaulay 模都是纯的 (第 2 节, 命题 2 及其推论)。

引理 1.-- 设 A 为 Noether 局部环，M 为有限生成的纯 A-模，x 为 \(\mathfrak{m}_{\mathrm{A}}\) 的一个元素。下列条件等价：

\begin{enumerate}
\def\labelenumi{(\roman{enumi})}
\item
  有 \(\dim(M/xM)=\dim(M)-1\)；
\item
  乘以 x 的映射 \(x_{M}\) 是单射。
\end{enumerate}

不妨设 M 非零。断言 (i) 等价于 x 不属于 Supp(M) 中满足 dim(A/p) = dim(M) 的任何极小元 p (VIII, § 3, 第 2 节, 命题 3)，而断言 (ii) 等价于 x 不属于 M 的任何伴随素理想 (IV, § 1, 第 1 节, 命题 2 的推论 2)。由于 M 是纯的，其伴随素理想就是 Supp(M) 的极小元，故 (i) 与 (ii) 等价。

设 A 为 Noether 局部环，M 为非零的有限生成 A-模。回忆（VIII, § 3, 第 2 节）：称 \((x_{1},\ldots,x_{r})\)（其元素属于 \(m_{A}\)）这样的序列对 M 是割的，如果有 \(\dim(M/(x_{1}M+\ldots+x_{r}M))=\dim(M)-r\)。

命题 4. 设 A 为 Noether 局部环，M 为非零的有限生成 A-模，\((x_{1},\ldots ,x_{r})\) 为 \(\mathfrak{m}_{A}\) 中元素的一个对 M 割的序列。下列条件等价：

\begin{enumerate}
\def\labelenumi{(\roman{enumi})}
\item
  A-模 M 是 Macaulay 的
\item
  序列 \((x_{1},\ldots,x_{r})\) 是 M-正则的，且 A-模 \(\mathrm{M}/(x_{1}\mathrm{M}+\ldots+x_{r}\mathrm{M})\) 是 Macaulay 的。
\end{enumerate}

设序列 \((x_{1},\ldots,x_{r})\) 是 M-正则的。于是有

\[
\dim (\mathrm{M}) = r + \dim (\mathrm{M} / (x _ {1} \mathrm{M} + \dots + x _ {r} \mathrm{M}))
\]

\[
\operatorname{prof} (\mathrm{M}) = r + \operatorname{prof} \left(\mathrm{M} / (x _ {1} \mathrm{M} + \dots + x _ {r} \mathrm{M})\right)
\]

(§ 1, 第 4 节, 命题 7)，由此得到蕴含关系 (ii)⇒(i)。

设 A-模 M 是 Macaulay 的，对 \(r\) 作归纳来证明 (ii)。当 \(r = 0\) 时断言是显然的。若 \(r \geqslant 1\)，则由归纳假设，A-模 \(N = M/(x_1M + \ldots + x_{r-1}M)\) 是 Macaulay 的，且有 \(\dim(N/x_rN) = \dim(N) - 1\)，因为序列 \((x_1, \ldots, x_r)\) 是割的；于是位似 \((x_r)_N\) 是单射的（引理 1），且 \(N/x_rN\) 是 Macaulay 的 (\(\mathrm{no.}\) 1, 例 3)，由此得到 (ii)。

定理 1.--- 设 \(A\) 为 Noether 局部环，M 为非零的有限生成 A-模，d 为 M 的维数，\(\mathbf{x} = (x_1, \ldots, x_d)\) 为 \(\mathfrak{m}_{\mathrm{A}}\) 中元素的一个对 M 割的序列，J 为它生成的理想。下列条件等价：

\begin{enumerate}
\def\labelenumi{(\roman{enumi})}
\item
  A-模 M 是 Macaulay 的
\item
  序列 x 是 M-正则的；
\item
  序列 x 对 M 是完全割的 (A, X, 第 157 页, 定义 2)；
\item
  M 关于理想 J 的重数 (VIII, § 7, 第 1 节, 定义 1) \(e_{\mathrm{J}}(\mathrm{M})\) 等于 A-模 M/JM 的长度；
\item
  对每个满足 \(1 \leqslant i \leqslant d\) 的整数 i，A-模 \(\mathrm{M}/(x_{1}\mathrm{M}+\ldots+x_{i-1}\mathrm{M})\) 都是纯的。
\end{enumerate}

(ii) 与 (iii) 的等价性由 A, X, 第 160 页, 定理 1 的推论 1 得出。由于 A-模 M/JM 具有有限长度 (VIII, § 3, 第 2 节, 定理 1)，(iii) 与 (iv) 的等价性由 VIII, § 4, 第 3 节, 命题 4 及第 4 节, 定理 3 得出。剩下要证明 (i)、(ii) 与 (v) 的等价性。

(i)⇒(v)：若 M 是 Macaulay 的，则诸模 \(\mathrm{M}/(x_{1}\mathrm{M}+\ldots+x_{i-1}\mathrm{M})\) 中的每一个都是 Macaulay 的 (命题 4)，从而是纯的。

(v)⇒(ii)：对诸模 M/(x₁M+\ldots+\(x_{i-1}\){}M) 中的每一个应用引理 1 即得。

(ii)⇒(i)：由命题 4 得出，因为 M/JM 具有有限长度，从而是 Macaulay 的。

\subsection{Macaulay 模的强割子集与商}

设 A 为 Noether 环，M 为有限生成 A-模，S 为 A 的一个子集。按照第八章的约定，我们用 SM 表示 M 的子模 \(\sum_{s\in S}sM\)，并用 G 表示 A 的由 S 生成的理想。

引理 2.--- 设 \(\overline{\mathfrak{G}}\) 为 \(\mathfrak{G}\) 在 A/Ann(M) 中的像。我们有

\[
\operatorname{ht} (\overline {{\mathfrak {G}}}) = \operatorname{codim} (\operatorname{Supp} (\mathrm{M} / \mathrm{SM}), \operatorname{Supp} (\mathrm{M})) .
\]

此外，当 SM ≠ M 时，我们有

\[
\operatorname{ht} (\overline {{{\mathfrak {G}}}}) \leqslant \operatorname{Card} (S).
\]

用 a 表示 M 的零化子。由 II, § 4, 第 4 节, 命题 18 的推论，A-模 M/SM 的支撑集是 V(S + a)。它在 Supp(M) 中的余维数因此等于 V(S + a) 在 V(a) 中的余维数，后者又等于 V((S + a)/a) 在 Spec(A/a) 中的余维数，即 G 的高度。

设 SM ≠ M；当 S 无限时不等式 ht(G) ≤ Card(S) 是显然的，而当 S 有限时它由 VIII, § 3, 第 3 节, 命题 4 b) 得出。

定义 2.--- 设 A 为 Noether 环，M 为有限生成 A-模，S 为 A 的一个有限子集。称 S 对 M 是强割的，如果有

\[
\operatorname{Card} (\mathrm{S}) \leqslant \operatorname{codim} (\operatorname{Supp} (\mathrm{M} / \mathrm{SM}), \operatorname{Supp} (\mathrm{M})) .
\]

注.- 1) A 的每个满足 SM = M 的有限子集 S 都对 M 是强割的。当 SM ≠ M 时，由引理 2，为使 S 对 M 是强割的，必须且只需 Card(S) = codim(Supp(M/SM), Supp(M))，或等价地 ht(G) = Card(S)。

\begin{enumerate}
\def\labelenumi{\arabic{enumi})}
\setcounter{enumi}{1}
\tightlist
\item
  若环 A 是局部环且模 M 非零，则 \(m_{A}\) 的每个对 M 强割的子集 S 都对 M 是割的。事实上，由于 A-模 M/SM 非零，我们有
\end{enumerate}

\[
\operatorname{Card} (S) \leqslant \operatorname{codim} (\operatorname{Supp} (M / S M), \operatorname{Supp} (M)) \leqslant \dim (M) - \dim (M / S M)
\]

(VIII, § 1, 第 2 节, 命题 3 a))，由此得到我们的断言。

命题 5.--- 设 A 为 Noether 环，M 为有限生成 A-模，S 为 A 的一个有限子集。下列条件等价：

\begin{enumerate}
\def\labelenumi{(\roman{enumi})}
\item
  A 的子集 S 对 M 是强割的；
\item
  对 Supp(M/SM) 的每个元素 p，典范映射 \(A \to A_{\mathfrak{p}}\) 诱导一个从 S 到 \(\mathfrak{p}A_{\mathfrak{p}}\) 中一个对 \(M_{\mathfrak{p}}\) 割的子集上的双射。
\item
  (i) \(\Rightarrow\) (ii)：设 \(\mathfrak{p} \in \operatorname{Supp}(\mathrm{M} / \mathrm{SM})\)，并设 \(S'\) 为 S 在 \(A_{\mathfrak{p}}\) 中的像。集合 \(S'\) 含于极大理想 \(\mathfrak{p} A_{\mathfrak{p}}\) 中，并且有
\end{enumerate}

\[
\dim (\mathrm{M} _ {\mathfrak {p}} / \mathrm{S} ^ {\prime} \mathrm{M} _ {\mathfrak {p}}) = \operatorname{codim} (\mathrm{V} (\mathfrak {p}), \operatorname{Supp} (\mathrm{M} / \mathrm{SM}))
\]

(VIII, § 1, 第 4 节, 命题 9)。不等式 Card(S) ≤ codim(Supp(M/SM), Supp(M)) 以及 VIII, § 1, 第 2 节, 命题 3 b) 蕴含关系式

\[
\operatorname{Card} (S) + \dim \left(M _ {\mathfrak {p}} / S ^ {\prime} M _ {\mathfrak {p}}\right) \leqslant \operatorname{codim} \left(V (\mathfrak {p}), \operatorname{Supp} (M)\right) = \dim \left(M _ {\mathfrak {p}}\right).
\]

另一方面，由于 \(M_{p}\) 非零，我们有 \(\dim(M_{\mathfrak{p}}) \leqslant \text{Card}(S') + \dim(M_{\mathfrak{p}}/S'M_{\mathfrak{p}})\) (VIII, § 3, 第 2 节, 公式 (8))。于是由不等式 \(\text{Card}(S') \leqslant \text{Card}(S)\) 得到条件 (ii)。

\begin{enumerate}
\def\labelenumi{(\roman{enumi})}
\setcounter{enumi}{1}
\tightlist
\item
  (ii) \(\Rightarrow\) (i)：不妨设 SM \(\neq\) M。若条件 (ii) 成立，则对 Supp(M/SM) 的每个素理想 \(\mathfrak{p}\)
\end{enumerate}

\[
\operatorname{Card} (S) = \operatorname{Card} \left(S ^ {\prime}\right) \leqslant \dim \left(M _ {\mathfrak {p}}\right) = \operatorname{codim} (V (\mathfrak {p}), \operatorname{Supp} (M)),
\]

对它取下确界即得 (i)。

推论.--- 设 A 为 Noether 环，M 为有限生成 A-模。每个 M-正则序列都对 M 是强割的。

设 x 为一个 M-正则序列，J 为它生成的 \(A\) 的理想。对每个素理想 \(\mathfrak{p} \in \operatorname{Supp}(\mathrm{M/JM})\)，x 在 \(A_{\mathfrak{p}}\) 中的像是 \(\mathrm{M}_{\mathfrak{p}}\)-正则的、由 \(\mathfrak{pA}_{\mathfrak{p}}\) 中元素组成的序列，因而是 \(\mathrm{M}_{\mathfrak{p}}\) 的割序列 (VIII, § 3, 第 2 节, 命题 3 的推论)。

命题 6. 设 A 为 Noether 环，M 为有限生成的 Macaulay A-模，S 为 A 的一个对 M 强割的有限子集。则 A-模 M/SM 是 Macaulay 的。

对每个极大理想 \(\mathfrak{m} \in \operatorname{Supp}(\mathrm{M} / \mathrm{SM})\)，S 在 \(A_{\mathfrak{m}}\) 中的像对 \(M_{\mathfrak{m}}\) 是割的 (命题 5)。由于 \(M_{\mathfrak{m}}\) 是 Macaulay \(A_{\mathfrak{m}}\)-模，\((\mathrm{M} / \mathrm{SM})_{\mathfrak{m}}\) 亦然 (命题 4)，由此得到本命题。

定理 2 (Macaulay-Cohen).--- 设 A 为 Noether 环，M 为有限生成 A-模。下列条件等价：

\begin{enumerate}
\def\labelenumi{(\roman{enumi})}
\item
  A-模 M 是 Macaulay 的
\item
  对 A 的每个由 A 中元素的 M-正则序列生成的理想 J，A-模 M/JM 都没有嵌入伴随素理想；
\item
  对 A 的每个对 M 强割的有限子集 S，A-模 M/SM 都没有嵌入伴随素理想。
\end{enumerate}

(i)⇒(iii)：设 S 为 A 的一个对 M 强割的有限子集。A-模 M/SM 是 Macaulay 的 (命题 6)，特别地没有嵌入伴随素理想 (第 2 节, 命题 2, a))。

(iii)⇒(ii)：这由命题 5 的推论得出。

\begin{enumerate}
\def\labelenumi{(\roman{enumi})}
\setcounter{enumi}{1}
\tightlist
\item
  (ii) → (i)：设 \(\mathfrak{p} \in \text{Supp}(M)\)；我们来证明 \(A_{\mathfrak{p}}\) -模 \(M_{\mathfrak{p}}\) 是 Macaulay 的。对整数 dim( \(M_{\mathfrak{p}}\) ) 作归纳。若 dim( \(M_{\mathfrak{p}}\) ) 为零，则 \(M_{\mathfrak{p}}\) 是有限长度的 \(A_{\mathfrak{p}}\) -模，从而是 Macaulay 的 (第 1 节, 例 1)。设 dim( \(M_{\mathfrak{p}}\) ) 非零，即 p 不是 Supp(M) 的极小元 (VIII, § 1, 第 4 节, 命题 9 a))。由于 M 没有嵌入伴随素理想，p 不含于 M 的任何伴随素理想之中，故存在 p 的元素 x，使乘以 x 的映射 \(x_{M}\) 是单射的 (§ 1, 注 2)。于是乘以 x 的映射 \(x_{M_{\mathfrak{p}}}\) 是单射的，且有 dim( \(M_{\mathfrak{p}}/xM_{\mathfrak{p}}\) ) \textless{} dim( \(M_{\mathfrak{p}}\) ) (VIII, § 3, 第 2 节, 命题 3)。对 A-模 M/xM 以及 Supp(M/xM) 的素理想 p 应用归纳假设，\(A_{\mathfrak{p}}\) -模 \(M_{\mathfrak{p}}/xM_{\mathfrak{p}}\) 是 Macaulay 的，这蕴含 \(A_{\mathfrak{p}}\) -模 \(M_{\mathfrak{p}}\) 是 Macaulay 的 (第 3 节, 命题 4)。
\end{enumerate}

\subsection{Macaulay 环}

定义 3.--- 称环 A 为 Macaulay 的，或称为 Macaulay 环，如果它是 Noether 的且 A-模 A 是 Macaulay 的。

例.- 1) 每个 Artin 环都是 Macaulay 环 (第 1 节, 例 1)。2) Macaulay 环没有嵌入伴随素理想 (第 2 节, 命题 2)。反之，设 A 为没有嵌入伴随素理想的维数 ≤ 1 的 Noether 环；对 A 的每个非空有限强割子集 S，A-模 A/SA 的维数 ≤ 0，从而是 Macaulay 的 (第 1 节, 例 1)；因此 A 是 Macaulay 环 (第 4 节, 定理 2)。特别地，维数 ≤ 1 的约化 Noether 环是 Macaulay 环。

\begin{enumerate}
\def\labelenumi{\arabic{enumi})}
\setcounter{enumi}{2}
\item
  维数 \(\leqslant 2\) 的正规 Noether 环是 Macaulay 环 (\(\S 1, n^{\circ} 10\), 定理 4 之前的正文)。反之，设 \(A\) 为一个 Macaulay 环，它在每个素理想处的局部环 \(A_{\mathfrak{p}}\)（素理想 \(\mathfrak{p}\) 的高度 \(\leqslant 1\)）都是整闭的；则 A 是正规的 (\(\S 1, n^{\circ} 10\), 参见定理 4)。
\item
  若 A 是 Macaulay 环，则对 A 的每个乘性子集 S，\(S^{-1}A\) 也是 Macaulay 环 (\(n^{\circ}\) 1, 注)。反之，若对 A 的每个极大理想 m，环 \(A_{m}\) 都是 Macaulay 环，则环 A 是 Macaulay 的 (\(n^{\circ}\) 1, 定义 1)。
\item
  设 A 为 Noether 环，J 为 \(A\) 的理想。为使 A/J 是 Macaulay 环，必须且只需它是 \(A\)-Macaulay 模 (第 1 节, 例 4)。
\item
  设 A 为 Noether 局部环，J 为由 A-正则序列生成的 A 的理想。为使 A/J 是 Macaulay 环，必须且只需 A 本身是 (第 3 节, 例 5 与命题 4)。
\item
  为使 Noether 局部环 A 是 Macaulay 的，必须且只需它具有一个由 A-正则序列生成的定义理想：这由第 3 节, 命题 3 以及下述事实得出：m\_A 中元素的 A-正则序列生成一个定义理想当且仅当其长度为 dim(A) (VIII, § 3, 第 2 节, 定理 1 及命题 3 的推论)。特别地，每个正则 Noether 局部环都是 Macaulay 环 (VIII, § 5, 第 2 节, 定理 1)。更一般地，正则 Noether 局部环 A 对由 A-正则序列生成的理想的商是 Macaulay 环 (例 6)。
\end{enumerate}

命题 7.-- 设 A 为 Noether 环。下列条件等价：

\begin{enumerate}
\def\labelenumi{(\roman{enumi})}
\item
  A 是 Macaulay 环；
\item
  对 Spec(A) 的每个闭子集 F，有 \(\mathrm{prof}_{\mathrm{F}}(\mathrm{A}) = \mathrm{codim}(\mathrm{F})\)；
\item
  A 的每个理想 J 都含有一个长度为 ht(J) 的 A-正则序列；
\end{enumerate}

(iii') A 的每个极大理想 m 都含有一个长度为 ht(m) 的 A-正则序列；

\begin{enumerate}
\def\labelenumi{(\roman{enumi})}
\setcounter{enumi}{3}
\tightlist
\item
  对 A 的每个理想 J，有 \(\mathrm{Ext}_A^i (\mathrm{A} / \mathrm{J},\mathrm{A}) = 0\) 对 \(i <   \mathrm{ht}(\mathrm{J})\) 成立
\end{enumerate}

(iv') 对 A 的每个极大理想 m，有 Ext \(_{A}^{i}\) (A/m, A) = 0 对 i \textless{} ht(m) 成立；

\begin{enumerate}
\def\labelenumi{(\alph{enumi})}
\setcounter{enumi}{21}
\tightlist
\item
  对 A 的每个素理想 p 以及 \(A_{p}\) 的每个由 \(A_{p}\) 的极大割序列生成的理想 J，有 \(e_{\mathrm{J}}(\mathrm{A}_{\mathfrak{p}})=\mathrm{long}(\mathrm{A}_{\mathfrak{p}}/\mathrm{JA}_{\mathfrak{p}})\)；
\end{enumerate}

(v') 对 A 的每个极大理想 m，存在 \(A_{m}\) 的一个理想 J，它由 \(A_{m}\) 的一个极大割序列生成，且满足 \(e_{\mathrm{J}}(A_{\mathrm{m}})=\mathrm{long}(\mathrm{A}_{\mathrm{m}}/\mathrm{JA}_{\mathrm{m}})\)；

\begin{enumerate}
\def\labelenumi{(\roman{enumi})}
\setcounter{enumi}{5}
\tightlist
\item
  (Macaulay-Cohen 判别法) 对 A 的每个使得由 S 生成的理想 G 具有高度 Card(S) 的有限子集 S，A-模 A/G 都没有嵌入伴随素理想。
\end{enumerate}

(i) 与 (ii) 的等价性由第 1 节, 命题 1 的推论得出。由 § 1, 第 4 节, 定理 2 以及深度的定义，条件 (iii) 与 (iv)（相应地 (iii') 与 (iv')）意味着对 A 的每个理想（相应地每个极大理想）J 有 \(\mathrm{prof}_{\mathrm{A}}(\mathrm{J};\mathrm{A})\geqslant \mathrm{ht}(\mathrm{J})\)。于是有

\[
(i) \Leftrightarrow (ii) \Rightarrow (iii) \Leftrightarrow (iv) \Rightarrow (iii') \Leftrightarrow (iv').
\]

但 (iv') 蕴含：对 A 的每个极大理想 m，\(\operatorname{Ext}_{A_{m}}^{i}(\kappa(m), A_{m}) = 0\) 对 \(i < \dim(A_{m})\) 成立，于是 \(\operatorname{prof}(A_{m}) \geqslant \dim(A_{m})\)，从而 A 是 Macaulay 环。

(i)、(v) 与 (v') 的等价性由第 3 节, 定理 1 得出，而 (i) 与 (vi) 的等价性由第 4 节, 定理 2 得出。

\subsection{Macaulay 模与有限代数}

注.--- 设 \(\rho: A \to B\) 为一个环同态，\(\mathfrak{p} \in \operatorname{Spec}(A)\)。用 \(\overline{B}\) 表示环 \(\kappa(\mathfrak{p}) \otimes_{A} B\)。它等同于 \(S^{-1}B/\mathfrak{p}(S^{-1}B)\)，其中 \(S\) 是形如 \(\rho(A - \mathfrak{p})\) 的 \(B\) 的乘性子集；因此 \(\overline{B}\) 的素理想就是诸理想 \(\mathfrak{q}\overline{B}\)，其中 \(\mathfrak{q}\) 是 \(B\) 的包含 \(\mathfrak{p}B\) 且不与 \(S\) 相交的素理想，换言之，即 \(B\) 的位于 \(\mathfrak{p}\) 之上的素理想。对这样的理想 \(\mathfrak{q}\)，有 \(S \subset B - \mathfrak{q}\)，故 \(\overline{B}\) 在 \(\mathfrak{q}\overline{B}\) 处的局部环等同于 \(B_{\mathfrak{q}}/pB_{\mathfrak{q}}\)，亦即 \(\kappa(\mathfrak{p}) \otimes_{A} B_{\mathfrak{q}}\)。

类似地，若 N 是一个 B-模，则 \(\overline{\mathrm{B}}_{\mathfrak{q}\overline{\mathrm{B}}}\) -模 \(\left(\kappa(\mathfrak{p})\otimes_{\mathrm{A}}\mathrm{N}\right)_{\mathfrak{q}\overline{\mathrm{B}}}\) 等同于 \(\kappa(\mathfrak{p})\otimes_{\mathrm{A}}\mathrm{N}_{\mathfrak{q}}\)。再设 B-模 N 是有限生成的；由 Nakayama 引理，条件 \(\kappa(\mathfrak{p})\otimes_{\mathrm{A}}\mathrm{N}_{\mathfrak{q}}=0\) 等价于 \(\mathrm{N}_{\mathfrak{q}}=0\)。于是 \(\overline{\mathrm{B}}\) -模 \(\kappa(\mathfrak{p})\otimes_{\mathrm{A}}\mathrm{N}\) 的支撑集由诸理想 \(\mathfrak{q}\overline{\mathrm{B}}\) 组成，其中 \(\mathfrak{q}\) 取遍 \(\operatorname{Supp}_{\mathrm{B}}(\mathrm{N})\) 中位于 \(\mathfrak{p}\) 之上的素理想。特别地，为使模 \(\kappa(\mathfrak{p})\otimes_{\mathrm{A}}\mathrm{N}\) 非零，必须且只需存在 \(\operatorname{Supp}_{\mathrm{B}}(\mathrm{N})\) 的一个位于 \(\mathfrak{p}\) 之上的素理想。

命题 8.--- 设 \(\rho: A \to B\) 为 Noether 环之间的同态，\(N\) 为一个 \(B\) -模，它作为 \(A\) -模是有限型的。为使 \(A\) -模 \(N\) 是 Macaulay 的，必须且只需 \(B\) -模 \(N\) 是 Macaulay 的，并且对每对 \((\mathfrak{n}, \mathfrak{n}')\)（\(\mathrm{Supp}_{\mathrm{B}}(N)\) 中满足 \(\rho^{-1}(\mathfrak{n}) = \rho^{-1}(\mathfrak{n}')\) 的极大理想），有 \(\dim_{B_{\mathfrak{n}}} (N_{\mathfrak{n}}) = \dim_{B_{\mathfrak{n}}'} (N_{\mathfrak{n}}')\)。

A-模 B/Ann \(_B(N)\) 同构于有限生成 A-模 \(\text{End}_A(N)\) 的一个子模，因此是有限型的。把 A 换成 A/Ann \(_A(N)\)，B 换成 B/Ann \(_B(N)\)，我们就归结到下述情形：\(\rho\) 是单射且使 B 成为有限 A-代数，并且有 \(\text{Supp}_A(N) = \text{Spec}(A)\) 及 \(\text{Supp}_B(N) = \text{Spec}(B)\)。映射 \(f: \text{Spec}(B) \to \text{Spec}(A)\)（由 \(\rho\) 诱导）这时是满射，且 B 的素理想 \(\mathfrak{q}\) 是极大的当且仅当 \(f(\mathfrak{q})\) 是 A 的极大理想 (V, § 2, 第 1 节, 定理 1 及命题 1)。

设 m 为 A 的一个极大理想。由上述注，环 \(B_{m}\) 的包含 \(mB_{m}\) 的素理想是形如 \(qB_{m}\) 的理想，其中 \(q \in \text{Spec}(B)\) 是 B 的（必然为极大的）理想，且满足 \(f(\mathfrak{q}) = \mathfrak{m}\)。我们有

\[
\operatorname{prof} _ {\mathrm{A} _ {\mathfrak {m}}} \left(\mathrm{N} _ {\mathfrak {m}}\right) = \operatorname{prof} _ {\mathrm{B} _ {\mathfrak {m}}} \left(\mathfrak {m B} _ {\mathfrak {m}}; \mathrm{N} _ {\mathfrak {m}}\right) = \inf _ {\mathfrak {q} \in f ^ {- 1} (\mathfrak {m})} \left(\operatorname{prof} _ {\mathrm{B} _ {\mathfrak {q}}} \left(\mathrm{N} _ {\mathfrak {q}}\right)\right)
\]

(§ 1, 第 3 节, 命题 4 及第 5 节, 命题 8)，以及

\[
\dim_ {A _ {\mathfrak {m}}} \left(N _ {\mathfrak {m}}\right) = \dim_ {B _ {\mathfrak {m}}} \left(N _ {\mathfrak {m}}\right) = \sup _ {\mathfrak {q} \in f ^ {- 1} (\mathfrak {m})} \left(\dim_ {B _ {\mathfrak {q}}} \left(N _ {\mathfrak {q}}\right)\right)
\]

(VIII, § 2, 第 3 节, 定理 1 及 § 1, 第 4 节, 命题 9)。由于有 \(\mathrm{prof}_{\mathcal{B}_{\mathfrak{q}}}(\mathbf{N}_{\mathfrak{q}}) \leqslant \dim_{\mathcal{B}_{\mathfrak{q}}}(\mathbf{N}_{\mathfrak{q}})\)（对每个 \(\mathfrak{q} \in f^{-1}(\mathfrak{m})\)），本命题由这些等式得出。

推论 1.--- 设 \(\rho: A \to B\) 为 Noether 环之间的同态。若 B 是有限 A-代数且是 Macaulay A-模，则它是 Macaulay 环。此外，若 \(\rho\) 是单射，则对 A 的每个理想 a 有 ht(aB) = ht(a)，且对 B 的每个理想 b 有 ht(b) = ht(ρ \(^{-1}\) (b))。

第一个断言由命题 8 得出。设 \(\rho\) 是单射。设 \(\mathfrak{a}\) 为 A 的理想。我们有 \(\mathrm{ht}(\mathfrak{a}) = \mathrm{prof}_{\mathrm{A}}(\mathfrak{a};\mathrm{B})\)，因为 A-模 B 是 Macaulay 的且其支撑集等于 \(\mathrm{Spec}(\mathrm{A})\) (\(n^{\circ}1\), 命题 1 的推论)；又 \(\mathrm{ht}(\mathfrak{a}\mathrm{B}) = \mathrm{prof}_{\mathrm{B}}(\mathfrak{a}\mathrm{B};\mathrm{B})\)（同上引文），且 \(\mathrm{prof}_{\mathrm{A}}(\mathfrak{a};\mathrm{B}) = \mathrm{prof}_{\mathrm{B}}(\mathfrak{a}\mathrm{B};\mathrm{B})\) (\(\S 1, n^{\circ}3\), 命题 4)，由此得 \(\mathrm{ht}(\mathfrak{a}\mathrm{B}) = \mathrm{ht}(\mathfrak{a})\)。设 \(\mathfrak{b}\) 为 B 的理想。由前所述，有 \(\mathrm{ht}(\boldsymbol{\rho}^{-1}(\mathfrak{b})) = \mathrm{ht}(\boldsymbol{\rho}^{-1}(\mathfrak{b})\mathrm{B})\)。但 \(\boldsymbol{\rho}^{-1}(\mathfrak{b})\mathrm{B}\) 含于 \(\mathfrak{b}\)，故其高度小于 \(\mathrm{ht}(\mathfrak{b})\)，而有 \(\mathrm{ht}(\mathfrak{b}) \leqslant \mathrm{ht}(\boldsymbol{\rho}^{-1}(\mathfrak{b}))\)（由 VIII, \(\S 2, n^{\circ}3\), 定理 1, b)）。

推论 2.--- 设 A 为整闭的 Noether 环，B 为包含 A 的环。假设 B 作为 A-模是有限型的无扭模。若 B 是 Macaulay 环，则 A-模 B 是 Macaulay 的。

事实上，B 中位于 A 的同一理想之上的两个素理想具有相同的高度 (VIII, § 2, 第 3 节, 定理 2)。因此可以对 N = B 应用命题 8。

推论 3.--- 设 A 为整闭环，K 为其分式域，L 为一个有限 K-代数，使得 {[}L : K{]} \(1_{\mathrm{A}}\) 在 A 中可逆，B 为 L 的一个子 A-代数，且在 A 上有限。

\begin{enumerate}
\def\labelenumi{\alph{enumi})}
\item
  B 的子 A-模 A1B 是一个直和因子。
\item
  对 A 的每个理想 J，有不等式 \(\operatorname{prof}_{\mathrm{A}}(\mathrm{J};\mathrm{A})\geqslant\operatorname{prof}_{\mathrm{B}}(\mathrm{JB};\mathrm{B})\)。
\item
  若 B 是 Macaulay 环，则 A 亦然。
\end{enumerate}

K-线性映射 \(\mathrm{Tr}_{\mathrm{L/K}}: \mathrm{L} \to \mathrm{K}\) 把 B 映到 A 中 (V, § 1, 第 6 节, 命题 17 的推论 2)，因而通过限制定义了一个 A-线性映射 \(t: \mathrm{B} \to \mathrm{A}\)。对每个 \(x \in \mathrm{A}\)，有 \(t(x1_{\mathrm{B}}) = [\mathrm{L}: \mathrm{K}]x\)，由此得到 a)。

由 § 1, 第 3 节, 命题 4，有 \(\operatorname{prof}_{\mathrm{A}}(\mathrm{J};\mathrm{B}) = \operatorname{prof}_{\mathrm{B}}(\mathrm{JB};\mathrm{B})\)；而由 a) 及 § 1, 第 1 节, 注 4，有 \(\operatorname{prof}_{\mathrm{A}}(\mathrm{J};\mathrm{A}) \geqslant \operatorname{prof}_{\mathrm{A}}(\mathrm{J};\mathrm{B})\)，由此得到 b)。

若环 B 是 Noether 的，则 A 亦然：事实上，由 a)，有 \(\mathfrak{a}\mathrm{B} \cap \mathrm{A} = \mathfrak{a}\)（对 A 的每个理想 \(\mathfrak{a}\)）；于是 A 的每个递增理想序列 \((\mathfrak{a}_n)_{n \in \mathbf{N}}\) 都是稳定的，因为序列 \((\mathfrak{a}_n\mathrm{B})_{n \in \mathbf{N}}\) 是稳定的。在 c) 的假设下，A-模 B 是 Macaulay 的（推论 2），从而 A-模 A 亦然 (第 1 节, 例 2)。

例.--- 推论 3 特别适用于下述两种情形：

\begin{enumerate}
\def\labelenumi{\alph{enumi})}
\item
  取一个 Noether 整闭环 A、其分式域的一个有限次 n 可分扩张 L（n 满足 \(n1_{A}\) 在 A 中可逆），并取 B 为 A 在 L 中的整闭包 (V, § 1, 第 6 节, 命题 18 的推论 1)。
\item
  取一个 Noether 整闭环 B 以及 B 的自同构的一个有限群 G，使得 \(\mathrm{Card(G)1_B}\) 在 B 中可逆。我们取 A 为 B 中在 G 的作用下不变的元素所组成的环。我们来验证这是 a) 的一个特殊情形。群 G 作用在 B 的分式域 L 上，而 L 对这个作用的不变元域就是 A 的分式域 K (V, § 1, 第 9 节, 命题 23 的推论)。L 是 K 的 Galois 扩张，从而更是可分扩张；其 Galois 群同构于 G，故 {[}L : K{]} 等于 Card G。{[}L : K{]} \(_{1_{B}}\) 的逆在 G 下不变，故 {[}L : K{]} \(_{1_{A}}\) 在 A 中可逆。由于 B 是整闭的，环 A（即 K ∩ B）是整闭的，且 B 是 A 在 L 中的整闭包（同上引文, 命题 22）。
\end{enumerate}

特别地，若 B 是 Macaulay 环，则 A 亦然。

\subsection{Macaulay 模与平坦代数}

命题 9.— 设 \(\rho: A \to B\) 为 Noether 环之间的同态，M 为有限生成 A-模，N 为有限生成 B-模且在 A 上平坦。用 \({}^{a}\rho: \operatorname{Spec}(B) \longrightarrow \operatorname{Spec}(A)\) 表示与 \(\rho\) 相伴的映射。下列条件等价：

\begin{enumerate}
\def\labelenumi{(\roman{enumi})}
\item
  B-模 M ⊗A N 是 Macaulay 的；
\item
  \((\kappa(\mathfrak{p})\otimes_{\mathrm{A}}\mathrm{B})\) -模 \(\kappa(\mathfrak{p})\otimes_{\mathrm{A}}\mathrm{N}\) 对每个 \(\mathfrak{p}\in\mathrm{Supp}_{\mathrm{A}}(\mathrm{M})\) 是 Macaulay 的，且 \(\mathrm{A}_{\mathfrak{p}}\) -模 \(\mathrm{M}_{\mathfrak{p}}\) 对每个 \(\mathfrak{p}\in{}^{a}\rho(\mathrm{Supp}_{\mathrm{B}}(\mathrm{N}))\) 是 Macaulay 的；
\item
  对 \(\mathrm{Supp}_{\mathrm{B}}(\mathrm{N})\) 的每个其在 A 中的原像 \(\mathfrak{p}\) 属于 \(\mathrm{Supp}_{\mathrm{A}}(\mathrm{M})\) 的极大理想，\(\mathrm{A}_{\mathfrak{p}}\) -模 \(\mathrm{M}_{\mathfrak{p}}\) 与 \((\kappa(\mathfrak{p}) \otimes_{\mathrm{A}} \mathrm{B})\) -模 \(\kappa(\mathfrak{p}) \otimes_{\mathrm{A}} \mathrm{N}\) 都是 Macaulay 的。
\end{enumerate}

此外，若 B-模 N 是忠实平坦的，则这些条件蕴含 A-模 M 是 Macaulay 的。

设 q 为 B 的一个属于 \(M \otimes_{A} N\) 的支撑集的素理想。令 \(\mathfrak{p} = \rho^{-1}(\mathfrak{q})\)。由于模 \((\mathrm{M} \otimes_{\mathrm{A}} \mathrm{N})_{\mathfrak{q}}\) 等同于 \(M_{p} \otimes_{A_{p}} N_{q}\)，模 \(M_{p}\) 与 \(N_{q}\) 都非零，\(\kappa(\mathfrak{p}) \otimes_{A} N_{\mathfrak{q}}\) 亦然 ( \(n^{\circ} 6\) , 注)。\(A_{p}\) -模 \(N_{q}\) 同构于 \(N_{p}\) 的一个分式模，故是平坦的，并且有等式

\[
\operatorname{prof} _ {\mathrm{B} _ {\mathfrak {q}}} \left(\left(\mathrm{M} \otimes_ {\mathrm{A}} \mathrm{N}\right) _ {\mathfrak {q}}\right) = \operatorname{prof} _ {A_ {\mathfrak {p}}} \left(\mathrm{M} _ {\mathfrak {p}}\right) + \operatorname{prof} _ {\mathrm{B} _ {\mathfrak {q}}} (\kappa (\mathfrak {p}) \otimes_ {\mathrm{A}} \mathrm{N} _ {\mathfrak {q}})
\]

\[
\dim_ {B _ {q}} \left(\left(M \otimes_ {A} N\right) _ {q}\right) = \dim_ {A _ {p}} \left(M _ {p}\right) + \dim_ {B _ {q}} \left(\kappa (p) \otimes_ {A} N _ {q}\right)
\]

(§ 1, 第 6 节, 命题 11, b) 及注)，其中每一项都是 ≥ 0 的整数。考虑到 B \(_{q}\) -模 κ(p) ⊗A N \(_{q}\) 是 Macaulay 的当且仅当它作为 (κ(p) ⊗A B \(_{q}\) )-模是 Macaulay 的 (第 1 节, 例 4)，由此推出下面两个条件等价：

(α) Bq-模 (M ⊗A N)q 是 Macaulay 的；

(β) \(A_{p}\) -模 \(M_{p}\) 与 \((\kappa(\mathfrak{p})\otimes_{A}B_{\mathfrak{q}})\) -模 \(\kappa(\mathfrak{p})\otimes_{A}N_{\mathfrak{q}}\) 都是 Macaulay 的。现在来证明 (iii) 蕴含 (i)。设 n 为 B 的一个属于 \(M\otimes_{A}N\) 的支撑集的极大理想，\(\mathfrak{p}=\rho^{-1}(\mathfrak{n})\)。由前所述，有 \(\mathfrak{p}\in\operatorname{Supp}_{A}(M)\cap{}^{a}\rho(\operatorname{Supp}_{B}(N))\)；条件 (iii) 与第 6 节的注蕴含上面的条件 (β) 在 q=n 时成立。于是 \(B_{n}\) -模 \((M\otimes_{A}N)_{n}\) 是 Macaulay 的，由此得到 (i)。

蕴含关系 (ii)⇒ (iii) 是显然的；我们来证明 (i) 蕴含 (ii)。设 B-模 M ⊗\_A N 是 Macaulay 的。设 p 为 Supp\_A(M) 的一个元素。不妨设 (κ(p) ⊗\_A B)-模 κ(p) ⊗\_A N 非零，即存在 Supp\_B(N) 的一个位于 p 之上的素理想 q (第 6 节, 注)。B\_q-模 (M ⊗\_A N)\_q 是 Macaulay 的 (第 1 节, 例 3)；由前面已证的蕴含关系 (α) ⇒ (β) 及第 6 节的注推出：A\_p-模 M\_p 与 (κ(p) ⊗\_A B)-模 (κ(p) ⊗\_A N) 都是 Macaulay 的，由此得到 (ii)。

此外，若 N 在 A 上忠实平坦，则有 \(\kappa(\mathfrak{p})\otimes_{\mathrm{A}}\mathrm{N}\neq 0\)（对每个 \(\mathfrak{p}\in\operatorname{Spec}(\mathrm{A})\)），于是 \({}^{a}\rho(\operatorname{Supp}_{\mathrm{B}}(\mathrm{N}))=\operatorname{Spec}(\mathrm{A})\) (第 6 节, 注)，从而 (ii) 蕴含 M 是 Macaulay 的。

推论 1.--- 设 \(\rho: A \to B\) 为 Noether 局部环之间的一个局部同态，M 为非零的有限生成 A-模，N 为非零的有限生成 B-模且为平坦 A-模。为使 B-模 M \(\otimes_{A}N\) 是 Macaulay 的，必须且只需 A-模 M 是 Macaulay 的且 B/ \(\mathfrak{m}_{A}\)B-模 N/ \(\mathfrak{m}_{A}\)N 是 Macaulay 的。

事实上，由于 \(N/m_{A}N\) 非零，N 是忠实平坦 A-模 (I, § 3, 第 1 节, 定义 1)。

推论 2.-- 设 \(\rho: A \to B\) 为 Noether 环之间的一个同态，它使 \(B\) 成为忠实平坦 \(A\)-模，并设 \(M\) 为一个有限型 \(A\)-模。为使 \(B\)-模 \(M \otimes_{A} B\) 是 Macaulay 的，必须且只需 \(A\)-模 \(M\) 是 Macaulay 的，且 \(\kappa(\mathfrak{p}) \otimes_{A} B\) 对每个 \(\mathfrak{p} \in \operatorname{Supp}(M)\) 都是 Macaulay 环。

推论 3.--- 设 A 为 Noether 环，B 为有限平坦 A-代数，M 为有限生成的 Macaulay A-模。则 B-模 M⊗ \(_{A}\) B 是 Macaulay 的。

事实上，对每个 \(\mathfrak{p} \in \operatorname{Spec}(A)\)，环 \(\kappa(\mathfrak{p}) \otimes_{A} B\) 是有限 \(\kappa(\mathfrak{p})\)-代数，从而是 Macaulay 的 (第 5 节, 例 1)，于是应用命题 9。

推论 4.--- 设 A 为 Noether 环，J 为 A 的理想，M 为有限生成 A-模。用 \(\hat{A}\) 与 \(\hat{M}\) 分别表示 A 与 M 关于 J-进拓扑的分离完备化，用 S 表示 A 的乘性子集 \(1 + J\)。考察下列条件：

\begin{enumerate}
\def\labelenumi{(\roman{enumi})}
\item
  A-模 M 是 Macaulay 的
\item
  \(\widehat{\mathrm{A}}\) -模 \(\widehat{\mathrm{M}}\) 是 Macaulay 的；
\item
  \(\mathrm{S}^{-1}\mathrm{A}\) -模 \(\mathrm{S}^{-1}\mathrm{M}\) 是 Macaulay 的；
\item
  对每个极大理想 \(\mathfrak{m} \in \operatorname{Supp}(\mathrm{M}) \cap \mathrm{V}(\mathrm{J})\)，\(\mathrm{A}_{\mathfrak{m}}\) -模 \(\mathrm{M}_{\mathfrak{m}}\) 是 Macaulay 的；
\item
  对每个素理想 \(\mathfrak{p} \in \operatorname{Supp}(M)\)（满足 \(\mathfrak{p} + J \neq A\)），\(A_{\mathfrak{p}}\) -模 \(M_{\mathfrak{p}}\) 是 Macaulay 的且环 \(\kappa(\mathfrak{p}) \otimes_{A} \widehat{A}\) 是 Macaulay 的。
\end{enumerate}

条件 (ii) 至 (v) 等价，且都被 (i) 蕴含。当理想 J 含于 A 的根时，条件 (i) 至 (v) 等价。

我们知道 (i) 蕴含 (iii) (第 1 节, 例 3)，而当 J 含于 A 的根时，(iii) 与 (i) 相同（因为 S 的元素这时都可逆）。

环 \(\widehat{\mathrm{A}}\) 是 Noether 的 (III, § 3, 第 4 节, 命题 8)；它等同于 S \(^{-1}\) A 关于 S \(^{-1}\) J-进拓扑的完备化，而 \(\widehat{\mathrm{A}}\) -模 \(\widehat{\mathrm{M}}\) 等同于 S \(^{-1}\) M 关于 S \(^{-1}\) J-进拓扑的完备化 (III, § 3, 第 5 节, 命题 12)。因此，为证明条件 (ii) 至 (v) 的等价性，可把 \(A\) 换成 S \(^{-1}\) A，J 换成 S \(^{-1}\) J，M 换成 S \(^{-1}\) M；换言之，可设 J 含于 A 的根。这时 A-模 \(\widehat{\mathrm{A}}\) 是忠实平坦的（同上引文, 命题 9）。

显然 (v) 蕴含 (iv)，且 (iv) 蕴含 (i)。

(i)⇒(ii)：设 m 为 A 的一个极大理想；则 m 是  的一个位于 m 之上的极大理想，且  的每个极大理想都可这样得到 (III, § 3, 第 4 节, 命题 8)。环 κ(m) ⊗A  是一个域，从而是 Macaulay 环；若 A-模 M 是 Macaulay 的，则由命题 9 的蕴含关系 (iii)⇒(i)，Â-模 M 亦然。

(ii)⇒(v)：\(\hat{A}\) -模 \(\hat{M}\) 同构于 \(M\otimes_{A}\hat{A}\) (III, § 3, 第 4 节, 定理 3)；若它是 Macaulay 的，则由命题 9 的 (i)⇒(ii) 推出：\(\kappa(\mathfrak{p})\otimes_{A}\hat{A}\) 对每个 \(\mathfrak{p}\in\mathrm{Supp}(M)\) 都是 Macaulay 环，且 A-模 M 是 Macaulay 的。

命题 10.- 设 \(\rho: A \to B\) 为 Noether 环之间的一个使 B 成为平坦 A-模的同态。下列条件等价：

\begin{enumerate}
\def\labelenumi{(\roman{enumi})}
\item
  B 是 Macaulay 环；
\item
  对 B 的每个素理想 q，环 \(A_{\rho^{-1}(\mathfrak{q})}\) 与 \(\kappa(\rho^{-1}(\mathfrak{q}))\otimes_{A}B\) 都是 Macaulay 的；
\item
  对 B 的每个极大理想 n，环 \(A_{\rho^{-1}(\mathfrak{n})}\) 与 \(\kappa(\rho^{-1}(\mathfrak{n})) \otimes_{A} B\) 都是 Macaulay 的。
\end{enumerate}

此外，若 B 在 A 上忠实平坦，则这些条件蕴含 A 是 Macaulay 环。

这是命题 9 当 M = A、N = B 时的特殊情形。

推论 1.--- 在 Macaulay 环上有限的平坦代数都是 Macaulay 环。

推论 2. 设 A 为 Macaulay 环，n 为正整数；则 \(\mathrm{A}[\mathrm{X}_1,\dots ,\mathrm{X}_n]\) 与 \(\mathrm{A}[[\mathrm{X}_1,\dots ,\mathrm{X}_n]]\) 都是 Macaulay 环。

只需处理 \(n = 1\) 的情形。环 A{[}T{]} 是 Noether 的 (A, VIII, § 1, 第 4 节, 推论 1)，且对每个域 \(k\)，环 \(k[T]\) 是 Macaulay 环 (第 5 节, 例 2)；因此环 A{[}T{]} 是 Macaulay 的 (命题 10)，而环 A{[}{[}T{]}{]} 是 Macaulay 的 (命题 9 的推论 4)。

推论 3.--- Noether Macaulay 环上的每个有限生成代数都是级链的。

事实上，这样的代数是 Macaulay 环上的多项式环的商，因而是 Macaulay 环的商（推论 2），从而是级链的 (第 2 节)。

\section{深度与同调维数}

\subsection{投射维数、内射维数、同调维数}

设 A 为环，M 为 A-模。回忆 (A, X, 第 134 页, 定义 1)，M 的投射维数记作 \(\mathrm{dp}_{\mathrm{A}}(\mathrm{M})\)，它是 M 的诸投射分解的长度组成的集合在 \(\overline{\mathbf{Z}}\) 中的下确界。我们有 \(\mathrm{dp}_{\mathrm{A}}(0) = -\infty\)，且当 M 非零时 \(\mathrm{dp}_{\mathrm{A}}(\mathrm{M}) \geqslant 0\)。为使 M 是投射的，必须且只需 \(\mathrm{dp}_{A}(\mathrm{M}) \leqslant 0\)。

例. 设 J 为 A 的一个由 A-正则序列 \(\mathbf{x} = (x_{1}, \ldots, x_{r})\) 生成的理想。我们有 dp \(_{A}\) (A/J) ≤ r。事实上，若 A 为零环这是显然的；否则，Koszul 复形 \(\mathbf{K}_{\bullet}(\mathbf{x}, \mathrm{A})\) 是 A/J 的一个长度为 r 的自由分解（同上引文, 第 159 页, 注 3）。此外，对每个 A-模 N，A-模 Ext \(_{A}^{r}\) (A/J, N) 与 N/JN 同构（同上引文）；因此，为使 dp \(_{A}\) (A/J) = r，必须且只需 J 不等于 A（同上引文, 第 134 页, 命题 1）。

类似地，M 的内射维数记作 \(\mathrm{di}_{\mathrm{A}}(\mathrm{M})\)，定义为 M 的诸内射分解的长度组成的集合的下确界。我们有 \(\mathrm{di}_{\mathrm{A}}(0) = -\infty\)，且 \(\mathrm{di}_{\mathrm{A}}(\mathrm{M}) \geqslant 0\)（若 \(M \neq 0\)）。为使 M 是内射的，必须且只需 \(\mathrm{di}_{\mathrm{A}}(\mathrm{M}) \leqslant 0\)。

命题 1.- 设 A 为环，M 为 A-模，n 为一个 \(\geqslant 0\) 的整数。下列条件等价：

\begin{enumerate}
\def\labelenumi{(\roman{enumi})}
\item
  有 \(\mathrm{di}_{\mathrm{A}}(\mathrm{M}) \leqslant n\)（换言之，M 具有一个长度 \(\leqslant n\) 的内射分解）；
\item
  对每个 A-模 N 及每个整数 \(i > n\)，有 \(\mathrm{Ext}_A^i (\mathrm{N},\mathrm{M}) = 0\)；
\item
  对 A 的每个理想 \(\mathfrak{a}\)，有 \(\operatorname{Ext}_{\mathrm{A}}^{n+1}(\mathrm{A}/\mathfrak{a},\mathrm{M})=0\)；
\item
  对 A-模的每个正合列
\end{enumerate}

\[
0 \rightarrow \mathrm{M} \longrightarrow \mathrm{I} ^ {0} \longrightarrow \mathrm{I} ^ {1} \longrightarrow \dots \longrightarrow \mathrm{I} ^ {n - 1} \longrightarrow \mathrm{Q} \rightarrow 0,
\]

其中诸 \(I^{i}\) 是内射的，A-模 Q 是内射的。

(i)⇒(ii)：这由 A, X, 第 100 页, 定理 1 得出。

(ii)⇒(iii)：这是显然的。

\((\mathrm{iii}) \Rightarrow (\mathrm{iv}) : \text{在 (iv) 的情形中，对每个 A-模 N，存在一个同构，它把 Ext}_A^1(\mathrm{N}, \mathrm{Q}) \text{映到 Ext}_A^{n+1}(\mathrm{N}, \mathrm{M}) (\mathrm{A}, \mathrm{X}, \text{p. 128, cor. 4}) ; \text{在假设 (iii) 下，A-模 Ext}_A^1(A / \mathfrak{a}, \mathrm{Q}) \text{ 对每个理想 } \mathfrak{a} \text{ （属于 A）均为零，且 Q 是内射的 (A, X, p. 93, prop. 11).}\)

(iv)⇒(i)：考察正合列 (A, X, 第 52 页)

\[
0 \longrightarrow \mathrm{M} \longrightarrow \mathrm{I} ^ {0} (\mathrm{M}) \longrightarrow \dots \longrightarrow \mathrm{I} ^ {n - 1} (\mathrm{M}) \longrightarrow \mathrm{K} ^ {n - 1} (\mathrm{M}) \longrightarrow 0;
\]

若条件 (iv) 成立，则 A-模 \(K^{n-1}(M)\) 是内射的，由此得到 (i)。

回忆 (A, X, 第 138 页, 定义 2)，环 A 的同调维数记作 dh(A)，它是 \(\overline{\mathbf{Z}}\) 中这样的整数 \(n\) 组成的集合的上确界：存在两个 A-模 M 与 N 使 Ext \(_{A}^{n}\) (M,N) 非零。它因此也是所有 A-模的投射（或内射）维数组成的集合的上确界；当 A 是 Noether 环时，可以只限于有限生成 A-模 (A, X, 第 139 页, 推论)。

\subsection{同调维数的局部化}

命题 2.--- 设 A 为环，M 与 N 为 A-模，i 为整数，S 为 \(A\) 的一个乘性子集。我们有 \(\mathrm{S}^{-1}\mathrm{A}\) -模的一个典范同构

\[
\mathrm{S} ^ {- 1} \operatorname{Tor} _ {i} ^ {A} (\mathrm{M}, \mathrm{N}) \longrightarrow \operatorname{Tor} _ {i} ^ {\mathrm{S} ^ {- 1} A} (\mathrm{S} ^ {- 1} \mathrm{M}, \mathrm{S} ^ {- 1} \mathrm{N}).
\]

若环 A 是 Noether 的且 A-模 M 有限生成，则有 S \(^{-1}\) A-模的一个典范同构

\[
\mathrm{S} ^ {- 1} \operatorname{Ext} _ {A} ^ {i} (\mathrm{M}, \mathrm{N}) \longrightarrow \operatorname{Ext} _ {\mathrm{S} ^ {- 1} \mathrm{A}} ^ {i} (\mathrm{S} ^ {- 1} \mathrm{M}, \mathrm{S} ^ {- 1} \mathrm{N}).
\]

由于 A-模 \(S^{-1}A\) 是平坦的，这由 A, X, 第 110 页, 命题 9 及第 111 页, 命题 10 得出。

推论.--- 设 A 为环，M 与 N 为 A-模，i 为整数。

\begin{enumerate}
\def\labelenumi{\alph{enumi})}
\item
  \(\mathrm{Tor}_{i}^{A}(\mathrm{M},\mathrm{N})\) 的支撑集含于 \(\mathrm{Supp}(\mathrm{M})\cap \mathrm{Supp}(\mathrm{N})\)，且当 A 是 Noether 的而 M 有限生成时，\(\mathrm{Ext}_{\mathrm{A}}^{i}(\mathrm{M},\mathrm{N})\) 的支撑集亦然。
\item
  设 A 是 Noether 的且模 M 与 N 都有限生成；若 A-模 M⊗A N 具有有限长度，则 Tor \(_{i}^{A}\) (M, N) 与 Ext \(_{A}^{i}\) (M, N) 亦然。
\end{enumerate}

若 p 是 A 的一个不属于 Supp(M) ∩ Supp(N) 的素理想，则模 \(M_{p}\) 或 \(N_{p}\) 之一为零，于是由命题 2 得到 a)。

为使 Noether 环上的有限生成模具有有限长度，必须且只需其支撑集由极大理想组成 (IV, § 2, 第 5 节, 命题 7)。在假设 b) 下，A-模 Tor \(_{i}^{A}\) (M, N) 与 Ext \(_{A}^{i}\) (M,N) 都是有限生成的 (A, X, 第 108 页, 推论)；因此断言 b) 由 a) 得出。

命题 3.--- 设 A 为 Noether 环，M 为有限生成 A-模，N 为 A-模。

\begin{enumerate}
\def\labelenumi{\alph{enumi})}
\tightlist
\item
  我们有
\end{enumerate}

\[
\mathrm{dp} _ {\mathrm{A}} (\mathrm{M}) = \sup _ {\mathfrak {p}} \mathrm{dp} _ {\mathrm{A} _ {\mathfrak {p}}} \left(\mathrm{M} _ {\mathfrak {p}}\right), \quad \mathrm{di} _ {\mathrm{A}} (\mathrm{N}) = \sup _ {\mathfrak {p}} \mathrm{di} _ {\mathrm{A} _ {\mathfrak {p}}} \left(\mathrm{N} _ {\mathfrak {p}}\right), \quad \mathrm{dh} (\mathrm{A}) = \sup _ {\mathfrak {p}} \mathrm{dh} \left(\mathrm{A} _ {\mathfrak {p}}\right),
\]

其中 p 取遍 A 的素理想（相应地极大理想）组成的集合。

\begin{enumerate}
\def\labelenumi{\alph{enumi})}
\setcounter{enumi}{1}
\tightlist
\item
  映射 \(\mathfrak{p} \mapsto \mathrm{dp}_{A_{\mathfrak{p}}}(\mathrm{M}_{\mathfrak{p}})\)（从 \(\operatorname{Spec}(A)\) 到 \(\overline{Z}\)）是上半连续的。
\end{enumerate}

我们来证明 a)。设 \(n\) 为一个 \(\geqslant 0\) 的整数。设 \(\mathrm{dp}_{A}(\mathbf{M}) < n\)。对 A 的每个素理想 \(\mathfrak{p}\) 及每个 \(\mathrm{A}_{\mathfrak{p}}\) -模 Q，\(\mathrm{A}_{\mathfrak{p}}\) -模 \(\mathrm{Ext}_{\mathrm{A}_{\mathfrak{p}}}^{n}(\mathrm{M}_{\mathfrak{p}},\mathrm{Q})\) 同构于 \((\mathrm{Ext}_{\mathrm{A}}^{n}(\mathrm{M},\mathrm{Q}))_{\mathfrak{p}}\) (命题 2)，故为零；由此推出不等式 \(\mathrm{dp}_{A_{\mathfrak{p}}}(\mathrm{M}_{\mathfrak{p}}) \leqslant \mathrm{dp}_{\mathrm{A}}(\mathrm{M})\) (A, X, 第 134 页, 命题 1)。反之，设有 \(\mathrm{dp}_{A_{\mathfrak{m}}}(\mathrm{M}_{\mathfrak{m}}) < n\)（对 A 的每个极大理想 \(\mathfrak{m}\)），并设 R 为一个 A-模。有 \((\operatorname{Ext}_{\mathrm{A}}^{n}(\mathrm{M}, \mathrm{R}))_{\mathfrak{m}} = 0\)（对每个 \(\mathfrak{m}\)）(命题 2)，故 \(\operatorname{Ext}_{\mathrm{A}}^{n}(\mathrm{M}, \mathrm{R}) = 0\) (II, § 3, 第 3 节, 定理 1 的推论 2)，这蕴含 \(\mathrm{dp}_{\mathrm{A}}(\mathrm{M}) < n\) (A, X, 第 134 页, 命题 1)。a) 的第一个等式由此得出。第二个等式同样证明，利用命题 1 (iii) 给出的内射维数的刻画。由于 dh(A) 是诸 A-模的内射维数组成的集合的上确界 (第 1 节)，第三个等式由此得出。

我们来证明 b)。设 \(\mathfrak{p}\) 为 A 的一个素理想，\(n = \mathrm{dp}_{\mathrm{A}_{\mathfrak{p}}}(M_{\mathfrak{p}})\)。我们来证明存在 \(\mathfrak{p}\) 在 \(\operatorname{Spec}(A)\) 中的一个邻域 U，使得有 \(\mathrm{dp}_{\mathrm{A}_{\mathfrak{q}}}(M_{\mathfrak{q}}) \leqslant n\)（对每个 \(\mathfrak{q} \in U\)）。当 \(n = +\infty\) 时这是显然的；当 \(n = -\infty\) 时，这由 M 的支撑集是闭集这一事实得出。现设 \(n\) 有限，并选取 A-模的一个正合列

\[
\mathrm{P} _ {n - 1} \stackrel {d _ {n - 1}} {\longrightarrow} \mathrm{P} _ {n - 2} \longrightarrow \dots \longrightarrow \mathrm{P} _ {0} \stackrel {d _ {0}} {\longrightarrow} \mathrm{M} \rightarrow 0,
\]

其中诸 \(P_{i}\) 是有限型自由模 (A, X, 第 53 页, 命题 6)。令 \(P = Ker d_{n-1}\)；它是一个有限表示模。\(A_{p}\) -模 \(P_{p}\) 是投射的 (A, X, 第 134 页, 命题 1)，因而是自由的 (II, § 3, 第 2 节, 命题 5 的推论 2)。由 II, § 5, 第 1 节, 命题 2 的推论，存在 \(A - p\) 的一个元素 f 使得 \(A_{f}\) -模 \(P_{f}\) 是自由的；于是对 Spec(A) 中由不含 f 的素理想组成的开集 U 的每个素理想 q，\(A_{q}\) -模 \(P_{q}\) 都是自由的。这就证明了 b)。

推论 1.--- 对 A 的每个乘性子集 S，我们有

\[
\mathrm{dp} _ {\mathrm{S} ^ {- 1} \mathrm{A}} (\mathrm{S} ^ {- 1} \mathrm{M}) \leqslant \mathrm{dp} _ {\mathrm{A}} (\mathrm{M}), \quad \mathrm{di} _ {\mathrm{S} ^ {- 1} \mathrm{A}} (\mathrm{S} ^ {- 1} \mathrm{N}) \leqslant \mathrm{di} _ {\mathrm{A}} (\mathrm{N}), \quad \mathrm{dh} (\mathrm{S} ^ {- 1} \mathrm{A}) \leqslant \mathrm{dh} (\mathrm{A}).
\]

推论 2.--- 若对 Supp(M) 的每个极大理想 m 有 dp \(_{A_m}\) (M \(_m\) ) \textless{} +∞，则 dp \(_A\) (M) \textless{} +∞。

事实上，由极大理想组成的 Supp(M) 的子空间 X 是拟紧的 (II, § 4, 第 2 节, 命题 8 与 9 及第 3 节, 命题 11 的推论 7)；映射 \(\mathfrak{m} \mapsto \mathrm{dp}_{\mathrm{A}_{\mathfrak{m}}}(\mathrm{M}_{\mathfrak{m}})\)（从 X 到 \(\overline{\mathbf{R}}\)）是上半连续的 (命题 3)，故有界 (TG, IV, 第 30 页, 定理 3 的推论)。

注. --- 设 A 为一个无限维的正则 Noether 环 (VIII, § 5, 习题 6 c))。下文将看到 (第 7 节, 定理 2 及 § 4, 第 1 节, 命题 1)，对 A 的每个极大理想 m 有 \(\mathrm{di}_{\mathrm{A_m}}(\mathrm{A_m}) = \mathrm{dh}(\mathrm{A_m}) < +\infty\)；于是命题 3 蕴含 \(\mathrm{di}_{\mathrm{A}}(\mathrm{A}) = \mathrm{dh}(\mathrm{A}) = +\infty\)。因此，函数 \(\mathfrak{p} \mapsto \mathrm{di}_{A_{\mathfrak{p}}}(\mathrm{N}_{\mathfrak{p}})\) 与 \(\mathfrak{p} \mapsto \mathrm{dh}(\mathrm{A}_{\mathfrak{p}})\) 一般而言不是上半连续的。

\subsection{Noether 环的同调维数}

设 A 为 Noether 局部环，M 为有限生成 A-模。回忆 (A, X, § 3, 第 6 节)，M 的一个分解

\[
\dots \longrightarrow \mathrm{L} _ {n} \stackrel {d _ {n}} {\longrightarrow} \mathrm{L} _ {n - 1} \longrightarrow \dots \longrightarrow \mathrm{L} _ {0} \stackrel {d _ {0}} {\longrightarrow} \mathrm{M} \rightarrow 0
\]

称为极小投射分解，如果每个模 \(L_{i}\) 都是有限型自由模，且复形 \(\kappa_{A} \otimes_{A} L\) 的微分为零。对每个整数 \(i \geqslant 0\)

我们有

\[
[ \mathrm{Ext} _ {\mathrm{A}} ^ {i} (\mathrm{M}, \kappa_ {\mathrm{A}}): \kappa_ {\mathrm{A}} ] = [ \mathrm{Tor} _ {i} ^ {\mathrm{A}} (\mathrm{M}, \kappa_ {\mathrm{A}}): \kappa_ {\Lambda} ] = \mathrm{rg} _ {\mathrm{A}} (\mathrm{L} _ {i})\tag{1}
\]

(A, X, 第 103 页, 例 3)。每个有限生成 A-模都具有这样一个分解 (A, X, 第 56 页, 命题 10)。

命题 4. 设 A 为 Noether 局部环，M 为有限生成 A-模，n 为一个 ≥ 0 的整数。下列条件等价：

\begin{enumerate}
\def\labelenumi{(\roman{enumi})}
\item
  有 \(\mathrm{dp}_{\mathrm{A}}(\mathrm{M}) < n\)
\item
  有 \(\mathrm{Tor}_n^{\mathrm{A}}(\mathrm{M},\kappa_{\mathrm{A}}) = 0\)；
\item
  有 \(\operatorname{Ext}_{\mathrm{A}}^{n}(\mathrm{M},\kappa_{\mathrm{A}}) = 0\)；
\item
  M 的每个极小投射分解的长度都 \textless{} n。
\end{enumerate}

蕴含关系 (i)⇒(ii) 与 (i)⇒(iii) 是直接的 (A, X, 第 100 页, 定理 1)。设 L 为 M 的一个极小投射分解；若 (ii) 或 (iii) 成立，则由 (1) 有 \(L_{n}=0\)。由于 M 的每个极小投射分解都同构于 L (A, X, 第 54 页, 命题 8)，我们得到 (iv)。蕴含关系 (iv)⇒(i) 是平凡的。

推论 1.--- 设 A 为 Noether 局部环，n 为一个 ≥ 0 的整数。下列条件等价：

\begin{enumerate}
\def\labelenumi{(\roman{enumi})}
\item
  有 \(\mathrm{dh}(\mathrm{A}) < n\)
\item
  有 \(\mathrm{Ext}_{A}^{i}(\mathrm{M},\mathrm{N}) = 0\) 且 \(\mathrm{Tor}_i^A (\mathrm{M},\mathrm{N}) = 0\)（对每对 A-模 \((\mathbf{M},\mathbf{N})\) 及每个整数 \(i\geqslant n\)）；
\item
  有 \(\mathrm{Tor}_n^{\mathrm{A}}(\kappa_{\mathrm{A}},\kappa_{\mathrm{A}}) = 0\)
\item
  有 \(\operatorname{Ext}_{\mathrm{A}}^{n}(\kappa_{\mathrm{A}},\kappa_{\mathrm{A}}) = 0\)；
\item
  有 \(\mathrm{dp}_{\mathrm{A}}(\kappa_{\mathrm{A}}) < n\)。
\end{enumerate}

显然 (i) 蕴含 (ii)，且 (ii) 蕴含 (iii) 与 (iv)。对 A-模 \(\kappa_{\mathrm{A}}\) 应用命题 4，可知 (iii) 与 (iv) 各自蕴含 (v)。我们证明 (v) 蕴含 (i)：若 \(\mathrm{dp}_{A}(\kappa_{\mathrm{A}}) < n\)，则对每个 A-模 M 有 \(\operatorname{Tor}_{n}^{\mathrm{A}}(\mathrm{M}, \kappa_{\mathrm{A}}) = 0\)；于是每个有限生成 A-模都具有有限投射维数 \(< n\) (命题 4)，这蕴含 \(\mathrm{dh}(\mathrm{A}) < n\) (A, X, 第 138 页, 命题 4)。

\[
\text { COROLLARY   2. } - \text { We   have   dh(A) = dp } _ {A} (\kappa_ {A})  .
\]

注.- 1) 设 A 为局部环。A-模 \(\mathrm{Tor}_{1}^{\mathrm{A}}(\kappa_{\mathrm{A}},\kappa_{\mathrm{A}})\) 同构于 \(\mathfrak{m}_{A}/\mathfrak{m}_{A}^{2}\)。因此，当 A 是 Noether 的时，为使 \(\mathrm{Tor}_{1}^{\mathrm{A}}(\kappa_{\mathrm{A}},\kappa_{\mathrm{A}})\) 为零，必须且只需 \(\mathfrak{m}_{\mathrm{A}}\) 为零，即 A 是一个域。于是推论 1 蕴含：同调维数为 0 的 Noether 局部环是域。

\begin{enumerate}
\def\labelenumi{\arabic{enumi})}
\setcounter{enumi}{1}
\tightlist
\item
  设 A 为 Noether 局部环，M 为具有有限投射维数 n 的有限生成 A-模，N 为非零的有限生成 A-模。A-模 Ext \(_{A}^{n}(M,N)\) 非零：事实上，设 L 为 M 的一个极小投射分解，d 为其微分。我们有正合列
\end{enumerate}

\[
\operatorname{Hom} _ {\mathrm{A}} \left(\mathrm{L} _ {n - 1}, \mathrm{N}\right) \xrightarrow {\operatorname{Hom} \left(d _ {n} , 1\right)} \operatorname{Hom} _ {\mathrm{A}} \left(\mathrm{L} _ {n}, \mathrm{N}\right) \longrightarrow \operatorname{Ext} _ {\mathrm{A}} ^ {n} (\mathrm{M}, \mathrm{N}) \rightarrow 0.
\]

由于 \(d_n \otimes 1_{\kappa_A}\) 为零，用 \(\kappa_A\) 作张量积得到一个同构 \(\kappa_A \otimes_A \operatorname{Hom}_A(L_n, N) \longrightarrow \kappa_A \otimes_A \operatorname{Ext}_A^n(M, N)\)，于是，考虑到上面的公式 (1)，

\[
\left[ \kappa_ {\mathrm{A}} \otimes_ {\mathrm{A}} \operatorname{Ext} _ {\mathrm{A}} ^ {n} (\mathrm{M}, \mathrm{N}): \kappa_ {\mathrm{A}} \right] = \left[ \operatorname{Ext} _ {\mathrm{A}} ^ {n} (\mathrm{M}, \kappa_ {\mathrm{A}}): \kappa_ {\mathrm{A}} \right] \left[ \kappa_ {\mathrm{A}} \otimes_ {\mathrm{A}} \mathrm{N}: \kappa_ {A} \right],
\]

由命题 4 与 Nakayama 引理，它非零。因此，M 的投射维数是使 Ext \(_{A}^{i}(M,N)\) 非零的最大整数 i。

\begin{enumerate}
\def\labelenumi{\arabic{enumi})}
\setcounter{enumi}{2}
\tightlist
\item
  设 A 为 Noether 环，M 为具有有限投射维数的有限生成 A-模，N 为支撑集等于 \(\operatorname{Spec}(A)\) 的有限生成 A-模。由前注及第 2 节的命题 2 与命题 3，投射维数 \(n\)（即 \(M\) 的）是这样的最大整数 \(i\)：\(\operatorname{Ext}_{A}^{i}(M,N)\) 非零；A-模 \(\operatorname{Ext}_{A}^{n}(M,N)\) 的支撑集是元素 \(\mathfrak{p}\)（属于 \(\operatorname{Spec}(A)\)）组成的集合，这里 \(\mathrm{dp}_{A_{\mathfrak{p}}}(M_{\mathfrak{p}})=n\)。
\end{enumerate}

可能存在这样的有限生成 A-模 M，其有限投射维数为 \(+\infty\)，且满足 \(\operatorname{Ext}_{\mathrm{A}}^{i}(\mathrm{M},\mathrm{A})=0\)（对充分大的 \(i\)）：*例如，A-模 \(\kappa_{\Lambda}\) 当 \(\Lambda\) 是非正则的 Gorenstein 局部环时即属此情形*。

命题 5.--- 设 A 为 Noether 环，M 为有限生成 A-模，n 为一个 ≥ 0 的整数。下列条件等价：

\begin{enumerate}
\def\labelenumi{(\roman{enumi})}
\item
  有 \(\mathrm{dp}_{\mathrm{A}}(\mathrm{M}) < n\)
\item
  对每个极大理想 \(\mathfrak{m}\)（属于 \(A\)），有 \(\operatorname{Ext}_{\mathrm{A}}^{n}(\mathrm{M},\mathrm{A}/\mathfrak{m}) = 0\)（相应地，有 \(\operatorname{Tor}_{n}^{\mathrm{A}}(\mathrm{M},\mathrm{A}/\mathfrak{m}) = 0\)）；
\item
  对 A 的每个极大理想 m，有 \(\operatorname{Ext}_{\mathrm{A}_{m}}^{n}(\mathrm{M}_{\mathfrak{m}},\mathrm{A}/\mathfrak{m})=0\)（相应地，有 \(\operatorname{Tor}_{n}^{\mathrm{A}_{\mathfrak{m}}}(\mathrm{M}_{\mathfrak{m}},A/\mathfrak{m})=0\)）。
\end{enumerate}

(i)⇒(ii)：这是显然的。

(ii)⇒(iii)：这由第 2 节的命题 2 得出。

\begin{enumerate}
\def\labelenumi{(\roman{enumi})}
\setcounter{enumi}{2}
\tightlist
\item
  (iii) \(\Rightarrow\) (i)：由命题 4，条件 (iii) 蕴含不等式 \(\mathrm{dp}_{A_{\mathfrak{m}}}(M_{\mathfrak{m}}) < n\) 对 A 的每个极大理想 \(\mathfrak{m}\) 成立，于是由命题 3 即得结论。
\end{enumerate}

注 4. — 设 A 为 Noether 环，\(n\) 为一个 \(\geqslant 0\) 的整数。若 \(\mathfrak{m}\) 与 \(\mathfrak{m}'\) 是 A 的两个不同的极大理想，则 A-模 \(\mathrm{Ext}_{\mathrm{A}}^{n}(\mathrm{A/m},\mathrm{A/m}')\) 与 \(\mathrm{Tor}_{n}^{\mathrm{A}}(\mathrm{A/m},\mathrm{A/m}')\) 都被 \(\mathfrak{m} + \mathfrak{m}'\) 零化，故均为零。用与命题 4 的推论 1 类似的证明，由命题 5 推出下列条件的等价性：

\begin{enumerate}
\def\labelenumi{(\roman{enumi})}
\item
  有 \(\mathrm{dh}(\mathrm{A}) < n\)
\item
  有 \(\operatorname{Ext}_A^i (\mathrm{M},\mathrm{N}) = 0\) 且 \(\operatorname{Tor}_i^{\mathrm{A}}(\mathrm{M},\mathrm{N}) = 0\)（对每对 A-模 \((\mathrm{M},\mathrm{N})\) 及每个整数 \(i\geqslant n\)）；
\item
  有 \(\mathrm{Tor}_n^{\mathrm{A}}(\mathrm{A} / \mathfrak{m},\mathrm{A} / \mathfrak{m}) = 0\)（对 A 的每个极大理想 \(\mathfrak{m}\)）；
\item
  有 \(\operatorname{Ext}_{\mathsf{A}}^{n}(\mathsf{A} / \mathfrak{m},\mathsf{A} / \mathfrak{m}) = 0\)（对 A 的每个极大理想 \(\mathfrak{m}\)）；
\item
  对 A 的每个极大理想 m，有 dp \(_{A}\) (A/m) \textless{} n。
\end{enumerate}

特别地，我们有 \(\mathrm{dh}(\Lambda)=\sup_{\mathfrak{m}}\mathrm{dp}_{\mathrm{A}}(\Lambda/\mathfrak{m})\)，其中 m 取遍 A 的极大理想组成的集合。

命题 6. 设 A 为 Noether 环，N 为 A-模，n 为一个 ≥ 0 的整数。下列条件等价：

\begin{enumerate}
\def\labelenumi{(\roman{enumi})}
\item
  有 \(\mathrm{di}_{\mathrm{A}}(\mathrm{N}) < n\)；
\item
  对每个素理想 \(\mathfrak{p}\)（属于 \(A\)），有 \(\mathrm{Ext}_{\mathrm{A}}^{n}(A/\mathfrak{p},\mathrm{N})=0\)；
\item
  对 A 的每个素理想 \(\mathfrak{p}\)，有 \(\operatorname{Ext}_{\mathrm{A}_{\mathfrak{p}}}^{n}(\kappa(\mathfrak{p}),\mathrm{N}_{\mathfrak{p}})=0\)。
\end{enumerate}

此外，若 A-模 N 有限生成，则这些条件等价于：

\begin{enumerate}
\def\labelenumi{(\roman{enumi})}
\setcounter{enumi}{3}
\tightlist
\item
  对 A 的每个极大理想 m，有 \(\operatorname{Ext}_{\mathrm{A}}^{i}(\mathrm{A}/\mathfrak{m},\mathrm{N})=0\) 对 \(n\leqslant i\leqslant n+\operatorname{ht}(\mathfrak{m})\) 成立。
\end{enumerate}

注意到条件 (iii) 等价于

(iii') 对 A 的每个素理想 \(\mathfrak{p}\)，有 \(\operatorname{Ext}_{\mathrm{A}}^{n}(\mathrm{A}/\mathfrak{p},\mathrm{N})\otimes_{\mathrm{A}}\kappa(\mathfrak{p})=0\)。

事实上，由于 \(\operatorname{Ext}_{\mathrm{A}}^{n}(\mathrm{A}/\mathfrak{p},\mathrm{N})\) 被 \(\mathfrak{p}\) 零化，A-模 \(\operatorname{Ext}_{\mathrm{A}}^{n}(\mathrm{A}/\mathfrak{p},\mathrm{N})\otimes_{\mathrm{A}}\kappa(\mathfrak{p})\) 同构于 \(\operatorname{Ext}_{\mathrm{A}}^{n}(\mathrm{A}/\mathfrak{p},\mathrm{N})\otimes_{\mathrm{A}}\mathrm{A}_{\mathfrak{p}}\)，从而同构于 \(\operatorname{Ext}_{\mathrm{A}_{\mathfrak{p}}}^{n}(\kappa(\mathfrak{p}),\mathrm{N}_{\mathfrak{p}})\) ( \(n^{\circ}\) 2, 命题 2)。

蕴含关系 (i)⇒(ii) 与 (i)⇒(iv) 由第 1 节的命题 1 得出，而蕴含关系 (ii)⇒(iii') 是显然的。

\begin{enumerate}
\def\labelenumi{(\roman{enumi})}
\setcounter{enumi}{2}
\tightlist
\item
  (iii)：对每个 A-模 M，令 T(M) = Ext \(_{A}^{n}\) (M,N)。设 (i) 不成立。则（第 1 节, 命题 1）存在 A 的理想 a 使 T(A/a) ≠ 0。设 p 为具有这一性质的理想中的极大者；我们来证明 p 是素理想。由 IV, § 1, 第 4 节, 定理 1 与定理 2，存在 A/p 的一个合成列 (M \(_{i}\) ) \(_{0\leqslant i\leqslant m}\)，使得每个商 M \(_{i}\) /M \(_{i+1}\) 都同构于某个模 A/p \(_{i}\)，其中 p \(_{i}\) (0 ≤ i ≤ m - 1) 是包含 p 的素理想。若对一切 i 都有 T(A/p \(_{i}\) ) 为零，则由正合列
\end{enumerate}

\[
\mathrm{T} \left(\mathrm{M} _ {0} / \mathrm{M} _ {i}\right) \longrightarrow \mathrm{T} \left(\mathrm{M} _ {0} / \mathrm{M} _ {i + 1}\right) \longrightarrow \mathrm{T} \left(\mathrm{M} _ {i} / \mathrm{M} _ {i + 1}\right)
\]

即可对 i 用归纳法推出 \(\mathrm{T}(\mathrm{M}_{0}/\mathrm{M}_{m})=\mathrm{T}(\mathrm{A}/\mathfrak{p})\) 为零，而事实并非如此。因此存在一个指标 i 使 \(\mathrm{T}(\mathrm{A}/\mathfrak{p}_{i})\neq0\)。由 p 的极大性，有 \(p=p_{i}\)，故 p 是素理想。

设 \(x\) 为 \(\mathrm{A} - \mathfrak{p}\) 的一个元素；由正合列

\[
0 \rightarrow \mathrm{A} / \mathfrak {p} \stackrel {{x _ {\mathrm{A} / \mathfrak {p}}}} {{\longrightarrow}} \mathrm{A} / \mathfrak {p} \longrightarrow \mathrm{A} / (\mathfrak {p} + x \mathrm{A}) \rightarrow 0
\]

得到正合列

\[
\mathrm{T} (\mathrm{A} / (\mathfrak {p} + x \mathrm{A})) \longrightarrow \mathrm{T} (\mathrm{A} / \mathfrak {p}) \stackrel {{u}} {{\longrightarrow}} \mathrm{T} (\mathrm{A} / \mathfrak {p}),
\]

其中 \(u = \operatorname{Ext}_{\mathrm{A}}^{n}(x_{\mathrm{A}}, 1_{\mathrm{N}}) - x_{\mathrm{T(A/p)}}\) (A, X, 第 89 页, 命题 6)。由 \(\mathfrak{p}\) 的极大性，有 \(\mathrm{T(A / (\mathfrak{p} + xA))} = 0\)，故同态 \(u\) 是单射。于是非零 A/p-模 \(\mathrm{T(A/p)}\) 是无扭的。这蕴含 \(\mathrm{T(A/p)} \otimes_{\mathrm{A}} \kappa(\mathfrak{p})\) 非零 (A, II, 第 117 页, 推论 1)，与 (iii') 矛盾。

设 A-模 N 有限生成，我们来证明 (iv) 蕴含 (iii)。设 p 为 A 的一个素理想，m 为包含 p 的一个极大理想。存在 A 的一个以 p 与 m 为端点的饱和素理想链。该链的长度 r 小于 ht(m)；因此在假设 (iv) 下，我们有

\[
\operatorname{Ext} _ {\mathrm{A} _ {\mathfrak {m}}} ^ {n + r} (\kappa (\mathfrak {m}), \mathrm{N} _ {\mathfrak {m}}) = \operatorname{Ext} _ {\mathrm{A}} ^ {n + r} (\mathrm{A} / \mathfrak {m}, \mathrm{N}) \otimes_ {A} \mathrm{A} _ {\mathfrak {m}} = 0.
\]

于是由 § 1, 第 7 节的引理 3 得出 \(\mathrm{Ext}_{\mathrm{A}_{\mathfrak{p}}}^{n}(\kappa(\mathfrak{p}),\mathrm{N}_{\mathfrak{p}})=0\)，这就证明了 (iii)。

注 5.— 设 N 为有限生成 A-模；条件「\(\operatorname{Ext}_{\mathrm{A}}^{n}(A/\mathfrak{m},\mathrm{N})=0\)（对 A 的每个极大理想 \(\mathfrak{m}\)）」并不必然蕴含 \(\operatorname{di}_{\mathrm{A}}(\mathrm{N})<n\)。例如，若环 A 是局部环且不是 Gorenstein 环 (第 7 节, 定义 1)，则有 \(\operatorname{Ext}_{A}^{n}(A/\mathfrak{m},\mathrm{A})=0\)（对 \(n<\operatorname{prof}(\mathrm{A})\)），但 \(\operatorname{di}_{\mathrm{A}}(\mathrm{A})=+\infty\)。

\subsection{对可消去元的商}

在本节中，A 表示一个环，x 是 A 的一个可消去元。

设 M 为 A-模，\(p: P \to M\) 为 M 的一个投射分解。由 A, X, 第 101 页, 定理 1 的推论，\(\mathrm{H}_{n}(\mathrm{P}/x\mathrm{P})\) 等同于 \(\mathrm{Tor}_{n}^{\mathrm{A}}(\mathrm{M}, \mathrm{A}/x\mathrm{A})\)，故当 n = 0 时同构于 M/xM，当 n = 1 时同构于 \(\mathrm{Ker}(x_{\mathrm{M}})\)，其余情形为零（同上引文, 第 102 页, 例 1）。考察 (A/xA)-模的复形 R 与 R'，使得

\[
\begin{array}{l l} \mathrm{R} _ {n} = \mathrm{P} _ {n} / x \mathrm{P} _ {n} & \quad \mathrm{R} _ {n} ^ {\prime} = 0 \qquad \text { for } n \geqslant 2, \\ \mathrm{R} _ {1} = \mathrm{Z} _ {1} (\mathrm{P} / x \mathrm{P}) & \quad \mathrm{R} _ {1} ^ {\prime} = (\mathrm{P} _ {1} / x \mathrm{P} _ {1}) / \mathrm{R} _ {1}, \\ \mathrm{R} _ {n} = 0 & \quad \mathrm{R} _ {n} ^ {\prime} = \mathrm{P} _ {n} / x \mathrm{P} _ {n} \qquad \text { for } n \leqslant 0, \end{array}
\]

其微分由 P 的微分诱导。复形 R(1) 与 R' 分别是 \(\mathrm{Ker}(x_{\mathrm{M}})\) 与 M/xM 的左分解，且有复形的正合列

\[
0 \rightarrow \mathrm{R} \longrightarrow \mathrm{P} / x \mathrm{P} \longrightarrow \mathrm{R} ^ {\prime} \longrightarrow 0.\tag{2}
\]

类似地，设 \(e: \mathrm{M} \to \mathrm{E}\) 为 M 的一个内射分解；用 K 表示复形 \(\operatorname{Ker}(x_{\mathrm{E}})\)。由 A, X, 第 101 页, 定理 1 的推论，\(\mathrm{H}_n(\mathrm{K})\) 等同于 \(\operatorname{Ext}_A^n (\mathrm{A} / x\mathrm{A},\mathrm{M})\)，故同构于 \(\operatorname{Ker}(x_{\mathrm{M}})\)（若 \(n = 0\)），同构于 \(\mathrm{M} / x\mathrm{M}\)（若 \(n = 1\)），其余情形为零（同上引文, 第 102 页, 例 1）。由 E 出发可导出 (A/xA)-模的复形 S 与 \(S'\)，使得

\[
\begin{array}{l l l} \mathrm{S} ^ {n} = \mathrm{K} ^ {n} & \quad \mathrm{S} ^ {\prime n} = 0 & \quad \text {for} n \leqslant 0, \\ \mathrm{S} ^ {1} = \mathrm{B} ^ {1} (\mathrm{K}) & \quad \mathrm{S} ^ {\prime 1} = \mathrm{K} ^ {1} / \mathrm{S} ^ {1}, \\ \mathrm{S} ^ {n} = 0 & \quad \mathrm{S} ^ {\prime n} = \mathrm{K} ^ {n} & \quad \text {for} n \geqslant 2. \end{array}
\]

复形 S 与 \(S'(-1)\) 分别是 \(\mathrm{Ker}(x_{\mathrm{M}})\) 与 M/xM 的右分解，且有复形的正合列

\[
0 \to \mathrm{S} \longrightarrow \operatorname{Ker} (x _ {\mathrm{E}}) \longrightarrow \mathrm{S} ^ {\prime} \to 0.\tag{3}
\]

设 N 为一个 (A/xA)-模，\(e': \mathrm{N} \to \mathrm{E}'\) 为 N 的一个内射分解。由正合列 (2) 可导出 (A/xA)-模的复形的正合列

\[
0 \to \operatorname{Homgr} _ {\mathrm{A} / x A} \left(\mathrm{R} ^ {\prime}, \mathrm{E} ^ {\prime}\right) \longrightarrow \operatorname{Homgr} _ {\mathrm{A} / x A} \left(\mathrm{P} / x \mathrm{P}, \mathrm{E} ^ {\prime}\right) \longrightarrow \operatorname{Homgr} _ {\mathrm{A} / x \mathrm{A}} \left(\mathrm{R}, \mathrm{E} ^ {\prime}\right)\rightarrow 0
\]

考察与这一正合列相伴的正合同调列。由 A, X, 第 100 页, 定理 1，对每个整数 \(n \geqslant 0\) 有同构

\[
\begin{array}{c} \mathrm{H} ^ {n} (\mathrm{Homgr} _ {A / x \mathrm{A}} (\mathrm{R} ^ {\prime}, \mathrm{E} ^ {\prime})) \longrightarrow \mathrm{Ext} _ {\mathrm{A} / x \mathrm{A}} ^ {n} (\mathrm{M} / x \mathrm{M}, \mathrm{N}) \\ \mathrm{H} ^ {n} (\mathrm{Homgr} _ {A / x A} (\mathrm{P} / x \mathrm{P}, \mathrm{E} ^ {\prime})) \longrightarrow \mathrm{H} ^ {n} (\mathrm{Homgr} _ {\mathrm{A}} (\mathrm{P}, \mathrm{E} ^ {\prime})) \longrightarrow \mathrm{Ext} _ {\mathrm{A}} ^ {n} (\mathrm{M}, \mathrm{N}) \\ \Pi^ {n} (\mathrm{Homgr} _ {\mathrm{A} / x \mathrm{A}} (\mathrm{R}, \mathrm{E} ^ {\prime})) = \mathrm{H} ^ {n - 1} (\mathrm{Homgr} _ {\mathrm{A} / x \mathrm{A}} (\mathrm{R} (1), \mathrm{E} ^ {\prime})) \longrightarrow \mathrm{Ext} _ {A / x A} ^ {n - 1} (\mathrm{Ker} (x _ {\mathrm{M}}), \mathrm{N}); \end{array}
\]

由此导出 (A/xA)-模的一个长正合列。

\[
\begin{array}{c} \dots \longrightarrow \operatorname{Ext} _ {\mathrm{A} / x \mathrm{A}} ^ {n} (\mathrm{M} / x \mathrm{M}, \mathrm{N}) \longrightarrow \operatorname{Ext} _ {\mathrm{A}} ^ {n} (\mathrm{M}, \mathrm{N}) \longrightarrow \operatorname{Ext} _ {\mathrm{A} / x \mathrm{A}} ^ {n - 1} (\operatorname{Ker} (x _ {\mathrm{M}}), \mathrm{N}) \\ \longrightarrow \operatorname{Ext} _ {\mathrm{A} / x \mathrm{A}} ^ {n + 1} (\mathrm{M} / x \mathrm{M}, \mathrm{N}) \longrightarrow \operatorname{Ext} _ {\mathrm{A}} ^ {n + 1} (\mathrm{M}, \mathrm{N}) \longrightarrow \dots \end{array}\tag{4}
\]

类似地，设 \(p' : \mathrm{P}' \to \mathrm{N}\) 为 \((A / x\mathrm{A})\) -模 \(\mathrm{N}\) 的一个投射分解。由正合列 (3)，我们导出 \((\mathrm{A} / x\mathrm{A})\) -模的复形的正合列

\[
0 \to \mathrm{Homgr} _ {\mathrm{A} / x \mathrm{A}} (\mathrm{P} ^ {\prime}, \mathrm{S}) \longrightarrow \mathrm{Homgr} _ {\mathrm{A} / x \mathrm{A}} (\mathrm{P} ^ {\prime}, \mathrm{Ker} (x _ {\mathrm{E}})) \longrightarrow \mathrm{Homgr} _ {\mathrm{A} / x \mathrm{A}} (\mathrm{P} ^ {\prime}, \mathrm{S} ^ {\prime}) \rightarrow 0
\]

考虑到同构 \(\mathrm{Homgr}_{\mathrm{A}/x_{\mathrm{A}}}(P',\mathrm{Ker}(x_{\mathrm{F}})) \longrightarrow \mathrm{Homgr}_{\mathrm{A}}(P',\mathrm{E})\)，相伴的正合同调列可写为

\[
\begin{array}{c} \dots \longrightarrow \operatorname{Ext} _ {\mathrm{A} / x A} ^ {n} (\mathrm{N}, \operatorname{Ker} (x _ {\mathrm{M}})) \longrightarrow \operatorname{Ext} _ {A} ^ {n} (\mathrm{N}, \mathrm{M}) \longrightarrow \operatorname{Ext} _ {\mathrm{A} / x A} ^ {n - 1} (\mathrm{N}, \mathrm{M} / x \mathrm{M}) \\ \longrightarrow \operatorname{Ext} _ {A / x A} ^ {n + 1} (\mathrm{N}, \operatorname{Ker} (x _ {\mathrm{M}})) \longrightarrow \operatorname{Ext} _ {A} ^ {n + 1} (\mathrm{N}, \mathrm{M}) \longrightarrow \dots \end{array}\tag{5}
\]

最后，考察 (A/xA)-模的复形的正合列

\[
0 \rightarrow \mathrm{R} \otimes_ {\mathrm{A} / x \mathrm{A}} \mathrm{P} ^ {\prime} \longrightarrow (\mathrm{P} / x \mathrm{P}) \otimes_ {\mathrm{A} / x \mathrm{A}} \mathrm{P} ^ {\prime} \longrightarrow \mathrm{R} ^ {\prime} \otimes_ {\mathrm{A} / x \mathrm{A}} \mathrm{P} ^ {\prime} \rightarrow 0
\]

它由正合列 (2) 导出；考虑到同构 \(P \otimes_{A} P' \longrightarrow (P/xP) \otimes_{A/xA} P'\)，相伴的正合同调列可写为

\[
\begin{array}{c} \ldots \longrightarrow \operatorname{Tor} _ {n} ^ {\mathrm{A}} (\mathrm{M}, \mathrm{N}) \longrightarrow \operatorname{Tor} _ {n} ^ {\mathrm{A} / x \mathrm{A}} (\mathrm{M} / x \mathrm{M}, \mathrm{N}) \longrightarrow \operatorname{Tor} _ {n - 2} ^ {\mathrm{A} / x \mathrm{A}} (\operatorname{Ker} x _ {\mathrm{M}}, \mathrm{N}) \\ \longrightarrow \operatorname{Tor} _ {n - 1} ^ {\mathrm{A}} (\mathrm{M}, \mathrm{N}) \longrightarrow \operatorname{Tor} _ {n - 1} ^ {\mathrm{A} / x \mathrm{A}} (\mathrm{M} / x \mathrm{M}, \mathrm{N}) \longrightarrow \ldots \end{array}\tag{6}
\]

命题 7.- 设 A 为环，\(x\) 为 A 的一个可消去元，M 为使乘以 x 的同态 \(x_{\mathrm{M}}\) 是单射的 A-模，N 为被 \(x\) 零化的 A-模，其中 n 为整数。(A/xA)-模的典范同态

\[
\begin{array}{r l r} \mathrm{Ext} _ {\mathrm{A} / x \mathrm{A}} ^ {n} (\mathrm{M} / x \mathrm{M}, \mathrm{N}) & \longrightarrow & \mathrm{Ext} _ {\mathrm{A}} ^ {n} (\mathrm{M}, \mathrm{N}), \\ \mathrm{Ext} _ {\mathrm{A}} ^ {n} (\mathrm{N}, \mathrm{M}) & \longrightarrow & \mathrm{Ext} _ {\mathrm{A} / x \mathrm{A}} ^ {n - 1} (\mathrm{N}, \mathrm{M} / x \mathrm{M}), \\ \mathrm{Tor} _ {n} ^ {A} (\mathrm{M}, \mathrm{N}) & \longrightarrow & \mathrm{Tor} _ {n} ^ {\mathrm{A} / x \mathrm{A}} (\mathrm{M} / x \mathrm{M}, \mathrm{N}) \end{array}
\]

由正合列 (4)、(5) 与 (6) 导出的上述同态都是同构。

推论.--- 设 A 为 Noether 局部环，M 为有限生成 A-模，J 为由序列 \((x_{1},\ldots,x_{r})\)（其元素属于 \(m_{A}\)，该序列既 A-正则又 M-正则）所生成的 A 的理想。我们有 dp \(_{A/J}\) (M/JM) = dp \(_{A}\) (M) 及 di \(_{A/J}\) (M/JM) = di \(_{A}\) (M) - r。

只需处理 \(r = 1\) 的情形。此时令 \(x = x_1\)。A-模 \(\mathrm{Ext}_{\mathrm{A} / x\mathrm{A}}^{n}(\mathrm{M} / x\mathrm{M},\kappa_{\mathrm{A}})\) 与 \(\mathrm{Ext}_{\mathrm{A}}^{n}(\mathrm{M},\kappa_{\mathrm{A}})\) 对每个整数 \(n\) 同构 (命题 7)；等式 \(\mathrm{dp}_{\mathrm{A} / x\mathrm{A}}(\mathrm{M} / x\mathrm{M}) = \mathrm{dp}_{\mathrm{A}}(\mathrm{M})\) 由第 3 节的命题 4 得出。类似地，A-模 \(\mathrm{Ext}_{A /x\mathrm{A}}^{n - 1}(\kappa_{\mathrm{A}},\mathrm{M} / x\mathrm{M})\) 与 \(\mathrm{Ext}_{\mathrm{A}}^{n}(\kappa_{\mathrm{A}},\mathrm{M})\) 对每个整数 \(n\) 同构，而等式 \(\mathrm{di}_{\mathrm{A} / x\mathrm{A}}(\mathrm{M} / x\mathrm{M}) = \mathrm{di}_{\mathrm{A}}(\mathrm{M}) - 1\) 由第 3 节的命题 6 得出。

\subsection{深度与投射维数}

定理 1 (Auslander-Buchsbaum).--- 设 A 为 Noether 局部环，M 为具有有限投射维数的有限生成 A-模。有等式

\[
\mathrm{dp} _ {\mathrm{A}} (\mathrm{M}) + \operatorname{prof} _ {\mathrm{A}} (\mathrm{M}) = \operatorname{prof} (\mathrm{A}).
\]

我们对 \(\mathrm{dp}_{\mathrm{A}}(\mathrm{M})\) 作归纳。

\begin{enumerate}
\def\labelenumi{\alph{enumi})}
\item
  若 dp \(_{A}\) (M) 为零，则 M 是非零的有限生成自由 A-模；其深度等于 prof(A) (§ 1, 第 1 节, 注 4)。
\item
  设 \(\mathrm{dp}_{A}(\mathbf{M}) = 1\)，我们选取 M 的一个极小投射分解
\end{enumerate}

\[
0 \to \mathrm{L} _ {1} \xrightarrow {d _ {1}} \mathrm{L} _ {0} \xrightarrow {d _ {0}} \mathrm{M} \to 0
\]

其中 A-模 L₀ 与 L₁ 是非零的有限生成自由模，故深度为 prof(A) (§ 1, 第 1 节, 注 4)。映射 \(1_{\kappa_{A}} \otimes d_{1} : \kappa_{A} \otimes_{A} L_{1} \longrightarrow \kappa_{A} \otimes_{A} L_{0}\) 为零，这蕴含 \(d_{1}\) 属于 \(\mathfrak{m}_{A} \operatorname{Hom}_{A}(L_{1}, L_{0})\)。由 § 1, 第 1 节的注 5，有 \(\operatorname{prof}_{A}(M) = \operatorname{prof}(A) - 1\)。

\begin{enumerate}
\def\labelenumi{\alph{enumi})}
\setcounter{enumi}{2}
\tightlist
\item
  设 \(\mathrm{dp}_{\mathrm{A}}(\mathrm{M}) > 1\)。选取一个正合列
\end{enumerate}

\[
0 \to \mathrm{N} \longrightarrow \mathrm{L} \longrightarrow \mathrm{M} \to 0
\]

其中 L 是有限生成自由 A-模。于是有 \(\operatorname{prof}_{\mathrm{A}}(\mathrm{L}) = \operatorname{prof}(\mathrm{A})\) ( \(\S 1\), 第 1 节, 注 4)，\(\mathrm{dp}_{\mathrm{A}}(\mathrm{N}) = \mathrm{dp}_{\mathrm{A}}(\mathrm{M}) - 1\) (A, X, 第 135 页, 推论 2 c))，故 \(\operatorname{prof}_{\mathrm{A}}(\mathrm{N}) = \operatorname{prof}(\mathrm{A}) - \mathrm{dp}_{\mathrm{A}}(\mathrm{N})\)（归纳假设），特别地 \(\operatorname{prof}_{\mathrm{A}}(\mathrm{N}) < \operatorname{prof}_{\mathrm{A}}(\mathrm{L})\)。于是由 \(\S 1\), 第 1 节的命题 1，得 \(\operatorname{prof}_{\mathrm{A}}(\mathrm{M}) = \operatorname{prof}_{\mathrm{A}}(\mathrm{N}) - 1\)，证明完毕。

注. 鉴于第 3 节命题 4 的推论 2，把定理 1 应用于 A-模 \(\kappa_{\mathrm{A}}\) 就推出：我们处于下述两种情形之一：

\begin{enumerate}
\def\labelenumi{(\roman{enumi})}
\item
  有 \(\mathrm{dp}_{\mathrm{A}}(\kappa_{\mathrm{A}}) = \mathrm{dh}(\mathrm{A}) = +\infty\)
\item
  有 \(\mathrm{dp}_{\mathrm{A}}(\kappa_{\mathrm{A}}) = \mathrm{dh}(\mathrm{A}) = \mathrm{prof}(\mathrm{A}) < +\infty\)。
\end{enumerate}

后文（§ 4, 第 2 节）将看到，(ii) 刻画了正则局部环。

推论 1.--- 保留定理 1 的假设。

\begin{enumerate}
\def\labelenumi{\alph{enumi})}
\item
  我们有 dp \(_{A}\) (M) ≤ prof(A)。为使等号成立，必须且只需极大理想 m \(_{A}\) 与 M 相伴。
\item
  我们有 \(\operatorname{prof}_A(M) \leqslant \operatorname{prof}(A)\)。为使等号成立，必须且只需 M 是自由的。
\item
  事实上，「 \(\operatorname{prof}_{A}(\mathsf{M}) = 0\) 」等价于「 \(\mathfrak{m}_A \in \operatorname{Ass}(A)\) 」(§ 1, 第 1 节, 注 2)。
\item
  事实上，「dp \(_{A}\) (M) = 0」等价于「M 是自由的」。
\end{enumerate}

推论 2. 保留定理 1 的假设，并设 A 还是 Macaulay 环。则 dp \(_{A}\) (M) 是两个正整数 dim(A) - dim \(_{A}\) (M) 与 dim \(_{A}\) (M) - prof(M) 之和。

特别地，这时 \(\mathrm{dp}_{\mathrm{A}}(\mathrm{M})\) 大于 \(\dim (\mathrm{A}) - \dim_{A}(\mathrm{M})\)，且当且仅当 M 是 Macaulay 的时取等号。

推论 3. 设 A 为 Noether 环，M 为具有有限投射维数的有限生成 A-模，i 为一个 \(\geqslant 0\) 的整数，N 为有限生成 A-模，F 为 A-模 \(\operatorname{Ext}_{\mathrm{A}}^{i}(\mathrm{M},\mathrm{N})\)（相应地 \(\operatorname{Tor}_{i}^{A}(\mathrm{M},\mathrm{N})\)）的支撑集。则我们有 \(\operatorname{prof}_{\mathrm{F}}(\mathrm{A}) \geqslant i\)。

事实上，设 \(\mathfrak{p} \in \mathrm{F}\)。由第 2 节的命题 2，有 \(\operatorname{Ext}_{\mathrm{A}_{\mathfrak{p}}}^{i}(\mathrm{M}_{\mathfrak{p}}, \mathrm{N}_{\mathfrak{p}}) \neq 0\)（相应地 \(\operatorname{Tor}_{i}^{\mathrm{A}_{\mathfrak{p}}}(\mathrm{M}_{\mathfrak{p}}, \mathrm{N}_{\mathfrak{p}}) \neq 0\)），故 \(i \leqslant \mathrm{dp}_{\mathrm{A}_{\mathfrak{p}}}(\mathrm{M}_{\mathfrak{p}}) \leqslant \mathrm{dp}_{\mathrm{A}}(\mathrm{M}) < +\infty\) (第 2 节, 命题 3)。定理 1 蕴含 \(\operatorname{prof}(\mathrm{A}_{\mathfrak{p}}) \geqslant i\)。因此 (§ 1, 第 5 节, 命题 8)

\[
\operatorname{prof} _ {F} (A) = \inf _ {\mathfrak {p} \in F} \operatorname{prof} (A _ {\mathfrak {p}}) \geqslant i.
\]

用 § 1, 第 5 节, 注 4 的术语，推论 3 的结论意味着模 \(\operatorname{Ext}_{A}^{i}(M,N)\) 与 \(\operatorname{Tor}_{i}^{A}(M,N)\) 的级 \(\geqslant i\)。它蕴含它们的支撑集在 \(\operatorname{Spec}(A)\) 中的余维数 \(\geqslant i\) ( \(\S 1\) , 第 7 节, 命题 12)。

推论 4.- 设 A 为 Noether Macaulay 环，M 为具有有限投射维数的有限生成 A-模。

\begin{enumerate}
\def\labelenumi{\alph{enumi})}
\tightlist
\item
  设 \(\mathfrak{p} \in \operatorname{Spec}(A)\)。用 \(\mathcal{C}(\mathfrak{p})\) 表示 \(\operatorname{Supp}(M)\) 的包含 \(\mathfrak{p}\) 的不可约分支组成的集合。我们有
\end{enumerate}

\[
\dim_ {A _ {\mathfrak {p}}} \left(M _ {\mathfrak {p}}\right) - \operatorname{prof} _ {A _ {\mathfrak {p}}} \left(M _ {\mathfrak {p}}\right) = \mathrm{dp} _ {A _ {\mathfrak {p}}} \left(M _ {\mathfrak {p}}\right) - \inf _ {X \in \mathscr {C} (\mathfrak {p})} \operatorname{codim} (X, \operatorname{Spec} (A)).
\]

\begin{enumerate}
\def\labelenumi{\alph{enumi})}
\setcounter{enumi}{1}
\item
  映射 \(\mathfrak{p}\mapsto\dim_{\mathrm{A}_{\mathfrak{p}}}(M_{\mathfrak{p}})-\operatorname{prof}_{\mathrm{A}_{\mathfrak{p}}}(M_{\mathfrak{p}})\)（从 Spec(A) 到 \(\overline{Z}\)）是上半连续的。
\item
  \(A\) 中使 A \(_{p}\) -模 M \(_{p}\) 为 Macaulay 的素理想 p 组成的集合在 \(\operatorname{Spec}(A)\) 中是开的且稠密。它与 \(\operatorname{Supp}(M)\) 的交在 \(\operatorname{Supp}(M)\) 中稠密。
\item
  可设 \(\mathfrak{p} \in \operatorname{Supp}(M)\)。令 \(\varphi(\mathfrak{p}) = \dim(A_{\mathfrak{p}}) - \dim_{A_{\mathfrak{p}}}(M_{\mathfrak{p}})\)。由上面的推论 2，我们有
\end{enumerate}

\[
\dim_ {A _ {\mathfrak {p}}} (M _ {\mathfrak {p}}) - \operatorname{prof} _ {A _ {\mathfrak {p}}} (M _ {\mathfrak {p}}) = \mathrm{dp} _ {A _ {\mathfrak {p}}} (M _ {\mathfrak {p}}) - \varphi (\mathfrak {p}).
\]

由于 A 是 Macaulay 环，我们有

\[
\dim (A _ {\mathfrak {p}}) = \operatorname{codim} (V (\mathfrak {p}), \operatorname{Spec} (A)) = \operatorname{codim} (V (\mathfrak {p}), X) + \operatorname{codim} (X, \operatorname{Spec} (A))
\]

对一切 \(X \in \mathcal{C}(\mathfrak{p})\) 成立 (§ 2, 第 2 节, 命题 2 b))；另一方面，我们有 (VIII, § 1, 第 4 节, 命题 9 及第 2 节, 注 3)

\[
\dim_ {A_ {\mathfrak {p}}} \left(\mathrm{M} _ {\mathfrak {p}}\right) = \operatorname{codim} (\mathrm{V} (\mathfrak {p}), \operatorname{Supp} (\mathrm{M})) = \sup _ {\mathrm{X} \in \mathscr {C} (\mathfrak {p})} \operatorname{codim} (\mathrm{V} (\mathfrak {p}), \mathrm{X})
\]

因而

\[
\varphi (\mathfrak {p}) = \inf _ {X \in \mathscr {C} (\mathfrak {p})} \operatorname{codim} (X, \operatorname{Spec} (A)).
\]

\begin{enumerate}
\def\labelenumi{\alph{enumi})}
\setcounter{enumi}{1}
\item
  设 \(\mathfrak{p} \in \operatorname{Spec}(A)\)，\(F\) 为 \(\operatorname{Supp}(M)\) 的不含 \(\mathfrak{p}\) 的诸不可约分支的并。对每个元素 \(\mathfrak{q}\)（属于 \(\operatorname{Spec}(A) - F\)），有 \(\mathcal{C}(\mathfrak{q}) \subset \mathcal{C}(\mathfrak{p})\)，于是由上面的公式得 \(\varphi(\mathfrak{q}) \geqslant \varphi(\mathfrak{p})\)。因此函数 \(\varphi\) 是下半连续的；断言 b) 于是由第 2 节的命题 3 得出。
\item
  设 U 为 Spec(A) 中使 \(M_{p}\) 为 Macaulay 的元素 p 组成的集合。条件 \(p \in U\) 等价于 \(\dim(M_{p}) - \text{prof}(M_{p}) \leqslant 0\)，故由 b) 知 U 是开的。由于 U 包含 \(\text{Spec}(A) - \text{Supp}(M)\)，只需证明 \(U \cap \text{Supp}(M)\) 在 \(\text{Supp}(M)\) 中稠密。对 \(\text{Supp}(M)\) 的每个极小素理想 p，\(A_{p}\) -模 \(M_{p}\) 具有有限长度 (IV, § 2, 第 5 节, 命题 7 的推论 2 及 § 1, 第 3 节, 命题 7 的推论 1)，故为 Macaulay 的；因此 U 与 \(\text{Supp}(M)\) 的每个不可约分支都相交。借助 II, § 4, 第 1 节的命题 1 即可完成证明。
\end{enumerate}

\subsection{深度与内射维数}

命题 8.--- 设 A 为 Noether 环，M 为有限生成 A-模。我们有 \(\dim_{A}(M) \leqslant \mathrm{di}_{A}(M)\)。

设 \(r\) 为一个小于 \(\dim_{\mathrm{A}}(\mathrm{M})\) 的正整数。存在素理想的一条饱和链 \(\mathfrak{p} \subset \mathfrak{p}_{1} \subset \ldots \subset \mathfrak{p}_{r-1} \subset \mathfrak{q}\)，使得 \(\mathfrak{p}\) 为 M 的支撑集的极小元；于是 \(\mathrm{A}_{\mathfrak{p}}\) -模 \(\mathrm{M}_{\mathfrak{p}}\) 具有有限长度，故有 \(\mathrm{Hom}_{\mathrm{A}_{\mathfrak{p}}}(\kappa(\mathfrak{p}), \mathrm{M}_{\mathfrak{p}}) \neq 0\)，从而 \(\mathrm{Ext}_{A_{\mathfrak{q}}}^{r}(\kappa(\mathfrak{q}), \mathrm{M}_{\mathfrak{q}}) \neq 0\) ( \(\S 1, n^{\circ} 7\) 引理 3)，这蕴含 \(\mathrm{di}_{\mathrm{A}}(\mathrm{M}) \geqslant r\) ( \(n^{\circ} 3\) , 命题 6)。

命题 9.--- 设 A 为 Noether 局部环，M 为具有有限内射维数的非零有限生成 A-模。则有 \(\mathrm{di}_{\mathrm{A}}(\mathrm{M}) = \mathrm{prof}(\mathrm{A})\)。

令 \(r = \mathrm{di}_{A}(\mathrm{M})\)。我们有 \(\operatorname{Ext}_{A}^{i}(\kappa_{A}, \mathrm{M}) = 0\) 对 \(i > r\) 成立，故有 \(\operatorname{Ext}_{\mathrm{A}}^{r}(\kappa_{\mathrm{A}}, \mathrm{M}) \neq 0\)（由 \(\mathfrak{n}^{\circ}\) 3 的命题 6, (iv)⇒(i)）。设 \(s\) 为 A 的深度，\((x_1, \ldots, x_s)\) 为 \(\mathfrak{m}_{\mathrm{A}}\) 中元素的一个 A-正则序列 (§ 1, \(\mathfrak{n}^{\circ}\) 4, 定理 2)；令 \(\mathrm{N} = \mathrm{A}/(x_1\mathrm{A} + \ldots + x_s\mathrm{A})\)。由 \(\mathfrak{n}^{\circ}\) 1 的例，我们有 \(\mathrm{dp}_{\mathrm{A}}(\mathrm{N}) = s\) 且 \(\operatorname{Ext}_{\mathrm{A}}^{s}(\mathrm{N}, \mathrm{M}) \neq 0\)，故 \(s \leqslant \mathrm{di}_{\mathrm{A}}(\mathrm{M}) = r\)。但 N 的深度为 0 (§ 1, \(\mathfrak{n}^{\circ}\) 4, 命题 7)，故存在 A-模的正合列

\[
0 \rightarrow \kappa_ {\mathrm{A}} \longrightarrow \mathrm{N} \longrightarrow \mathrm{N} ^ {\prime} \rightarrow 0.
\]

由此导出扩张模的一个正合列

\[
\operatorname{Ext} _ {\mathrm{A}} ^ {r} (\mathrm{N}, \mathrm{M}) \longrightarrow \operatorname{Ext} _ {\mathrm{A}} ^ {r} (\kappa_ {A}, \mathrm{M}) \longrightarrow \operatorname{Ext} _ {\mathrm{A}} ^ {r + 1} (\mathrm{N} ^ {\prime}, \mathrm{M});
\]

由于 \(\mathrm{Ext}_{\mathrm{A}}^{r + 1}(\mathrm{N}^{\prime},\mathrm{M}) = 0\) 且 \(\mathrm{Ext}_{\mathrm{A}}^{r}(\mathbf{k}_{\mathrm{A}},\mathrm{M})\neq 0\)，我们得到 \(\mathrm{Ext}_{\mathrm{A}}^{r}(\mathrm{N},\mathrm{M})\neq 0\)，故 \(r\leqslant \mathrm{dp}_{\mathrm{A}}(\mathrm{N}) = s\)。最终有 \(r = s\)，证明完毕。
